\subsection{幂零李群}

命题 12. 设 G 为有限维李群。L(G) 为幂零的充要条件是 G 具有一个幂零开子群。

设 \(\mathbf{G}\) 具有一个幂零开子群 \(\mathbf{G}_0\)。由命题 4 和命题 6 的第 2 号推论，\(\mathcal{C}^i\mathrm{L}(G_0) = \{0\}\)（当 \(i\) 足够大时）。因此 \(\mathrm{L}(G_0) = \mathrm{L}(G)\) 为幂零的。

设 \(\mathrm{L}\langle \mathrm{G}\rangle\) 为幂零的。若 \(\mathbf{K} = \mathbf{R}\) 或 \(\mathbf{C}\)，则由命题 4 的第 2 号推论，\(\mathbf{G}_0\) 是 \(\mathbf{G}\) 的单位元分支并且是幂零的，且 \(\mathbf{G}_0\) 在 \(\mathbf{G}\) 中是开的。若 \(\mathbf{K}\) 为超度量的，则命题 6 的第 2 号推论表明，存在一个开子群 \(\mathbf{G}_1\)，它是 \(\mathbf{G}\) 的子群，还存在一个整数 \(i > 0\) 和一个邻域 \(\mathbf{V}\)，其中包含 \(e\)，且位于 \(\mathbf{G}\) 中，使得 \(\mathrm{C}^i\mathrm{G}_1\cap \mathbf{V} = \{\varepsilon \}\)。于是，若 \(\mathbf{G}_0\) 是 \(\mathbf{G}_1\) 的充分小的子群，则 \(\mathrm{C}^i\mathrm{G}_0\subset \mathrm{V}\)，从而 \(\mathrm{C}^i\mathrm{G}_0 = \{\varepsilon \}\)，且 \(\mathbf{G}_0\) 为幂零的。

设 \(\mathfrak{g}\) 为幂零李代数。对应于 \(\mathfrak{g}\) 的 Hausdorff 级数 H(X, Y) 只有有限个非零项；我们知道（第二章，第 6 节，第 5 号，注 3），合成律 \((x, y) \mapsto H(x, y)\) 在 \(\mathfrak{g}\) 上定义了一个群结构。再设 \(\mathfrak{g}\) 完备且可赋范。显然，律 H 是连续多项式（《可微与解析流形》，R，附录）。因此，赋予律 H 的 \(\mathfrak{g}\) 是一个李群 G，称为与 \(\mathfrak{g}\) 相伴的李群。由第 4 节第 2 号的引理 2 和 3，\(\mathrm{L}(\mathrm{G}) = \mathfrak{g}\)。\(\phi\) 是从 \(\mathfrak{g}\) 到 \(\mathrm{G}\) 的恒等映射，是 \(\mathrm{G}\) 的一个指数映射，满足 \(\phi(\lambda x)\phi(\lambda' x) = \phi((\lambda + \lambda') x)\) 对所有 \(x \in \mathfrak{g}\)、\(\lambda \in \mathrm{K}\)、\(\lambda' \in \mathrm{K}\) 成立。每个李子代数 \(\mathfrak{h}\)（它是 \(\mathfrak{g}\) 的）具有拓扑补空间，都是 \(\mathrm{H}\) 的一个李子群 \(\mathrm{G}\)，且 \(\mathrm{L}(\mathrm{H}) = \mathfrak{h}\)。

命题 13. 设 G 是 R 或 C 上的有限维单连通幂零李群。

\begin{enumerate}
\def\labelenumi{(\roman{enumi})}
\item
  \(\exp_{\mathbf{G}}\) 是与 \(\mathrm{L}(\mathrm{G})\) 相伴的李群到 \(\mathbf{G}\) 上的一个同构。
\item
  \(\mathbf{G}\) 的每个积分子群都是 \(\mathbf{G}\) 的单连通李子群。
\end{enumerate}

设 \(\mathfrak{g} = \mathrm{L}(\mathrm{G})\)，它是幂零的（命题 12）。由于具有同一李代数的两个 \(\mathbf{R}\) 或 \(\mathbf{C}\) 上的单连通李群是同构的（第 6 节第 3 号定理 3(ii)），只需在 \(\mathrm{G}\) 是与 \(\mathfrak{g}\) 相伴的群时证明本命题。于是 (i) 和 (ii) 由命题之前所述内容推出。

命题 14. 设 G 是 R 或 C 上的有限维连通李群。

\begin{enumerate}
\def\labelenumi{(\roman{enumi})}
\item
  若 G 为幂零的，则 \(\exp_{G}\) 是 étale 的满射。
\item
  若 K = C 且 \(\exp_{G}\) 是 étale 的，则 G 为幂零的。
\end{enumerate}

设 \(G'\) 为 G 的万有覆叠空间。设 \(\phi\) 为从 \(G'\) 到 G 的典范态射。则 \(\exp_G = \phi \circ \exp_{G'}\)（第 6 节第 4 号命题 10），故 (i) 由命题 13(i) 得出。

若 \(\mathbf{K} = \mathbf{C}\) 且 exp 是 étale 的，则对所有 \(x \in \mathrm{L}(\mathrm{G})\)，\(\operatorname{ad} x\) 都没有属于 \(2i\pi (\mathbf{Z} - \{0\})\) 的特征值（第 6 节第 4 号命题 12 的推论）。对 \(\lambda x\) 应用这一点，其中 \(\lambda\) 遍历 \(\mathbf{C}\)，可得 ad \(x\) 的所有特征值都为零，因而 ad \(x\) 是幂零的。因此 \(\mathrm{L}(\mathrm{G})\) 是幂零的（第一章第 4 节定理 1 的推论 1），从而 G 是幂零的（命题 12）。

命题 15. 设 \(\mathbf{G}\) 是 \(\mathbf{R}\) 或 \(\mathbf{C}\) 上的有限维连通幂零李群，\(\mathbf{A}\) 是 \(\mathbf{G}\) 的一个积分子群。则 \(Z_{\mathrm{G}}(\mathbf{A})\) 是 \(\mathbf{G}\) 中李代数为 \(\mathfrak{z}_{\mathfrak{g}}(\mathrm{L}(\mathbf{A}))\) 的连通李子群。

由第 3 号命题 9，只需证明 \(Z_{\mathrm{G}}(\mathbf{A})\) 是连通的。设 \(g \in Z_{\mathrm{G}}(\mathbf{A})\)。存在 \(x \in \mathrm{L}(\mathbf{G})\)，使得 \(g = \exp x\)（命题 14）。于是 \(\operatorname{Ad} g|\mathrm{L}(\mathbf{A}) = 1\)（第 3 号命题 9），因而 \(\operatorname{Ad} g^n |\mathrm{L}(\mathbf{A}) = 1\)，对所有 \(n \in \mathbf{Z}\) 成立，从而 \(\exp (\operatorname{ad} nx)|\mathrm{L}(\mathbf{A}) = 1\)，对所有 \(n \in \mathbf{Z}\) 成立。由于映射

\[
\lambda \mapsto \exp (\operatorname{ad} \lambda x) | \mathrm{L} (\mathrm{A})
\]

从 K 到 \(\mathcal{L}(\mathrm{L}(\mathrm{A}),\mathrm{L}(\mathrm{G}))\) 的映射是多项式，所以 \(\exp(\operatorname{ad}\lambda x)|\mathrm{L}(\mathrm{A})=1\) 对所有 \(\lambda\in K\) 成立，即 \(\exp(\lambda x)\in Z_{\mathrm{G}}(\mathrm{A})\) 对所有 \(\lambda\in K\) 成立。

命题 16. 设 G 是 R 或 C 上的有限维幂零李群，A 是 G 中异于 G 的一个积分子群。则 \(\mathrm{N}_{\mathrm{G}}(\mathrm{A})\) 是 G 中异于 A 的连通李子群。

\(N_{\mathrm{G}}(\mathsf{A}) \neq \mathsf{A}\)（《代数》第一章第 6 节命题 8 的推论 1）。由第 4 号命题 11，只需证明 \(N_{\mathrm{G}}(\mathsf{A})\) 是连通的。设 \(g \in N_{\mathrm{G}}(\mathsf{A})\)。存在 \(x \in \mathrm{L}(\mathrm{G})\)，使得 \(g = \exp x\)（命题 14）。设 \(\mathrm{E}\) 是 \(\mathcal{L}(\mathrm{L}(\mathrm{G}))\) 的向量子空间，由 \(u \in \mathcal{L}(\mathrm{L}(\mathrm{G}))\) 组成，且 \(u(\mathrm{L}(\mathsf{A})) \subset \mathrm{L}(\mathsf{A})\)。于是 \(\operatorname{Ad} g^n \in \mathbb{E}\)，从而 \(\exp (\operatorname{ad} nx) \in \mathbb{E}\) 对所有 \(n \in \mathbf{Z}\) 成立。因此 \(\exp (\operatorname{ad} \lambda x) \in \mathbb{E}\) 对所有 \(\lambda \in K\) 成立，即 \(\exp (\lambda x) \in N_{\mathrm{G}}(\mathsf{A})\) 对所有 \(\lambda \in K\) 成立。

命题 17. 设 \(\mathfrak{g}\) 是 \(\mathbf{K}\) 上的有限维幂零李代数，\((\mathfrak{g}_0, \mathfrak{g}_1, \ldots, \mathfrak{g}_n)\) 是 \(\mathfrak{g}\) 的递减理想序列，满足 \(\mathfrak{g}_0 = \mathfrak{g}, \mathfrak{g}_n = \{0\}\)，且 \([\mathfrak{g}, \mathfrak{g}_i] \subset \mathfrak{g}_{i+1}\) 对 \(0 \leqslant i < n\)。设 \(\mathfrak{a}_1, \mathfrak{a}_2, \ldots, \mathfrak{a}_p\) 是 \(\mathfrak{g}\) 的向量子空间，并且每个 \(\mathfrak{g}_i\) 都是它与各个 \(\mathfrak{a}_j\) 的交的直和。赋予 \(\mathfrak{g}\) Hausdorff 合成律 H。设 \(\phi\) 是映射

\[
\left(x _ {1}, x _ {2}, \dots , x _ {p}\right) \mapsto x _ {1} \mathsf {H} x _ {2} \mathsf {H} \dots \mathsf {H} x _ {p}
\]

\[
o f a _ {1} \times a _ {2} \times \dots \times a _ {p}
\]

\begin{enumerate}
\def\labelenumi{(\roman{enumi})}
\item
  \(\phi\) 是从 \(a_{1} \times a_{2} \times \cdots \times a_{p}\) 到 \(g\) 的双射；
\item
  \(\phi\) 和 \(\phi^{-1}\) 都是多项式映射；
\item
  映射 \((x,y)\mapsto\phi^{-1}(\phi(x)\cdot\phi(y)^{-1})\)，从 \((\mathfrak{a}_{1}\times\mathfrak{a}_{2}\times\cdots\times\mathfrak{a}_{p})^{2}\) 到 \(a_{1}\times a_{2}\times\cdots\times a_{p}\)，是多项式。
\end{enumerate}

当 \(\dim \mathfrak{g} = 0\) 时本命题显然成立。设 \(\dim \mathfrak{g} > 0\)，并且本命题已对维数 \(<\dim \mathfrak{g}\) 成立。可以假定 \(\mathfrak{g}_{n-1} \neq \{0\}\)，于是 \(\mathfrak{g}_{n-1}\) 是 \(\mathfrak{g}\) 的非零中心理想。存在指标 \(j\)，使得 \(\mathfrak{b} = \mathfrak{g}_{n-1} \cap \mathfrak{a}_j \neq \{0\}\)。令 \(\mathfrak{g}' = \mathfrak{g}/\mathfrak{b}\)，\(\theta\) 为从 \(\mathfrak{g}\) 到 \(\mathfrak{g}'\) 的典范态射，\(\mathfrak{g}'_i = \theta(\mathfrak{g}_i)\) 且 \(\mathfrak{a}_i' = \theta(\mathfrak{a}_i)\)。则 \((\mathfrak{g}_0', \mathfrak{g}_1', \ldots, \mathfrak{g}_n')\) 是 \(\mathfrak{g}'\) 的递减理想序列，满足 \(\mathfrak{g}_0' = \mathfrak{g}'\)，\(\mathfrak{g}_n' = \{0\}\)，

\[
[ \mathfrak {g} ^ {\prime}, \mathfrak {g} _ {i} ^ {\prime} ] \subset \mathfrak {g} _ {i + 1} ^ {\prime}
\]

当 \(0 \leqslant i < n\) 时，且每个 \(\mathfrak{g}_i'\) 都是它与各个 \(\mathfrak{a}_j'\) 的交的直和。设 \(\phi'\) 是映射

\[
\left(x _ {1} ^ {\prime}, x _ {2} ^ {\prime}, \dots , x _ {p} ^ {\prime}\right) \mapsto x _ {1} ^ {\prime} \mathsf {H} x _ {2} ^ {\prime} \mathsf {H} \dots \mathsf {H} x _ {p} ^ {\prime}
\]

从 \(a_{1}^{\prime} \times a_{2}^{\prime} \times \cdots \times a_{p}^{\prime}\) 到 \(g^{\prime}\) 的映射。由归纳假设，\(\phi^{\prime}\) 是双射，且 \(\phi^{\prime}, \phi^{\prime-1}\) 是多项式映射。

设 \(x \in g\)。记

\[
\phi^ {\prime - 1} (\theta (x)) = \left(x _ {1} ^ {\prime} (x), x _ {2} ^ {\prime} (x), \dots , x _ {p} ^ {\prime} (x)\right).\tag{1}
\]

于是

\[
\theta (x) = x _ {1} ^ {\prime} (x) \mathsf {H} x _ {2} ^ {\prime} (x) \mathsf {H} \dots \mathsf {H} x _ {p} ^ {\prime} (x).\tag{2}
\]

设 \(\mathfrak{h}_1\) 是 \(\mathfrak{b}\) 在 \(\mathfrak{g}\) 中的一个向量补空间，它等于各个 \(a_k\)（当 \(k \neq j\) 时）之和与 \(\mathfrak{b}\) 在 \(a_j\) 中的一个补空间之和。存在从 \(\eta\) 到 \(\mathfrak{g}'\) 的双射 \(\mathfrak{h}_1\)，满足 \(\theta \circ \eta = \mathrm{Id}_{\mathfrak{g}'}\)。对所有 \(x \in \mathfrak{g}\)，记

\[
\zeta (x) = \eta (x _ {1} ^ {\prime} (x)) \mathsf {H} \eta (x _ {2} ^ {\prime} (x)) \mathsf {H} \dots \mathsf {H} \eta (x _ {p} ^ {\prime} (x)) \in \mathfrak {g}.\tag{3}
\]

\[
y (x) = \zeta (x) ^ {- 1} \text { H } x = (- \zeta (x))\cdot x.\tag{4}
\]

由 (2) 和 (3)，\(\theta (\zeta (x)) = \theta (x)\)，从而 \(y(x)\in \mathfrak{h}\)。最后记

\[
\psi (x) = \left(\eta \left(x _ {1} ^ {\prime} (x)\right), \dots , \eta \left(x _ {j} ^ {\prime} (x)\right) + y (x), \dots , \eta \left(x _ {p} ^ {\prime} (x)\right)\right) \in a _ {1} \times \dots \times a _ {p}.\tag{5}
\]

由于 \(y(x)\) 在 \(\mathfrak{g}\) 的中心中，

\[
\begin{array}{l l} \phi (\psi (x)) = \eta (x _ {1} ^ {\prime} (x)) \texttt {H} \dots \texttt {H} \eta (x _ {j} ^ {\prime} (x)) \texttt {H} \dots \texttt {H} \eta (x _ {p} ^ {\prime} (x)) \texttt {H} y (x) \\ = \zeta (x) \texttt {H} y (x) & \text {by (3)} \\ = x & \text {by (4)}. \end{array}
\]

因此 \(\phi \circ \psi = \mathrm{Id}_{\mathfrak{g}'}\)。现在设 \((x_{1}, x_{2}, \ldots, x_{p}) \in \mathfrak{a}_{1} \times \mathfrak{a}_{2} \times \dots \times \mathfrak{a}_{p}\)，并记 \(x = \phi(x_{1}, x_{2}, \ldots, x_{p}) = x_{1} \mathsf{H} x_{2} \mathsf{H} \dots \mathsf{H} x_{p}\)。于是

\[
\theta (x) = \theta (x _ {1}) \mathsf {H} \theta (x _ {2}) \mathsf {H} \dots \mathsf {H} \theta (x _ {p}),
\]

所以 \(x_{i}^{\prime}(x) = \theta (x_{i})\) 对所有 \(1\leqslant i\leqslant p\) 成立，从而

\[
\begin{array}{c} \zeta (x) = x _ {1} \mathsf {H} x _ {2} \mathsf {H} \dots \mathsf {H} (\eta \theta (x _ {j})) \mathsf {H} \dots \mathsf {H} x _ {p} \\ y (x) = x _ {j} - \eta \theta (x _ {j}). \end{array}
\]

由 (5)，

\[
\psi (x) = \left(x _ {1}, \dots , \eta \theta \left(x _ {j}\right) + x _ {j} - \eta \theta \left(x _ {j}\right), \dots , x _ {p}\right) = \left(x _ {1}, x _ {2}, \dots , x _ {p}\right).
\]

因此 \(\psi \circ \phi = \mathrm{Id}_{\mathfrak{a}_1 \times \ldots \times \mathfrak{a}_p}\)。这就证明了 (i)。由于 Hausdorff 律是多项式的，\(\phi\) 是多项式映射。由归纳假设，\(\phi'^{-1}\) 是多项式映射；根据公式 (1)，函数 \(x_j'\) 是多项式的，因而 \(\zeta\) 是多项式映射（公式 (3)），\(y\) 是多项式的（公式 (4)），而 \(\psi\) 是多项式映射（公式 (5)）。这就证明了 (ii)。(iii) 由 (i)、(ii) 以及 Hausdorff 律为多项式这一事实推出。

幂零李群的例子。设 G 是 \(\mathbf{GL}(n,\mathbf{K})\) 的严格下三角子群。它是 \(\mathbf{GL}(n,\mathbf{K})\) 的李子群，且 \(\mathrm{L}(\mathrm{G})\subset\mathrm{gl}(n,\mathbf{K})\) 是对角线为零的下三角矩阵李代数（第 3 节第 10 号命题 36）。由第二章第 4 节第 6 号注，G 是幂零的。以下设 K=R 或 C。由于 G 同胚于 \(\mathbf{K}^{n(n-1)/2}\)，G 是单连通的。L(G) 到 G 的指数映射就是映射

\[
u \mapsto \exp u = \sum_ {k \geqslant 0} \frac {u ^ {k}}{k !} = \sum_ {k = 0} ^ {n - 1} \frac {u ^ {k}}{k !}
\]

（第 6 节第 4 号例）。由命题 13，指数映射是从流形 L(G) 到流形 G 的同构。第 6 节第 9 号命题 17 给出其逆双射 log。在 \(\mathbf{K}^{n}\) 上赋予一个范数。由《谱理论》第一章第 4 节第 9 号，当 g ∈ G 且 \textbar\textbar g − 1\textbar\textbar{} \textless{} 1 时，

\[
\log g = \sum_ {k \geqslant 1} \frac {(- 1) ^ {k - 1}}{k} (g - 1) ^ {k}
\]

即

\[
\log g = \sum_ {k = 1} ^ {n - 1} \frac {(- 1) ^ {k - 1}}{k} (g - 1) ^ {k}.\tag{6}
\]

但 (6) 的两边都是关于 \(g\) 的解析函数，对 \(g \in \mathbf{G}\) 成立，因而对所有 \(g \in \mathbf{G}\) 都相等。

命题 18. 设 \(k\) 是交换域。V 是维数 \(>0\) 的 \(k\) 上向量空间，且维数大于 0，G 是 \(\mathbf{GL}(\mathbf{V})\) 的一个子群，其元素都是幺幂的。

\begin{enumerate}
\def\labelenumi{(\roman{enumi})}
\item
  存在一个非零元素 \(v\)，属于 \(V\)，使得 \(gv = v\) 对所有 \(g \in G\) 成立。
\item
  存在 V 的一个基 B，使得对所有 \(g \in G\)，\(g\) 关于 B 的矩阵都是下三角的，且其所有对角元都等于 1。
\item
  群 G 是幂零的。
\end{enumerate}

\begin{enumerate}
\def\labelenumi{(\alph{enumi})}
\tightlist
\item
  首先设 \(k\) 是代数闭域，且 \(G\) 的恒等表示是单的。设 \(a, b\) 属于 \(G\)。于是
\end{enumerate}

\[
\operatorname{Tr} (a (b - 1)) = \operatorname{Tr} (a b - 1) - \operatorname{Tr} (a - 1) = 0 - 0 = 0
\]

因为 \(ab - 1\) 和 \(a - 1\) 都是幂零的。由于由 \(\mathcal{L}(\mathrm{V})\) 中的 \(\mathbf{G}\) 生成的向量子空间就是 \(\mathcal{L}(\mathrm{V})\)（《代数》第八章第 4 节命题 2 的推论 1），\(\operatorname{Tr}(u(b - 1)) = 0\) 对所有 \(u \in \mathcal{L}(\mathrm{V})\) 成立，从而 \(b = 1\)。因此 \(\mathrm{G} = \{1\}\)。

\begin{enumerate}
\def\labelenumi{(\alph{enumi})}
\setcounter{enumi}{1}
\item
  现在转到一般情形。设 \(\overline{k}\) 是 \(k\) 的代数闭包，\(\overline{\mathrm{V}} = \mathrm{V} \otimes_k \overline{k}\)，且 \(\overline{\mathbf{G}} \subset \mathbf{GL}(\overline{\mathrm{V}})\) 是由 \(a \otimes 1\)（\(a \in \mathbf{G}\)）组成的集合。令 \(\mathrm{W}\)（分别地，\(\mathrm{W}'\)）为 \(\mathrm{V}\)（分别地，\(\overline{\mathrm{V}}\)）中在 \(\mathrm{G}\)（分别地，\(\overline{\mathbf{G}}\)）下不变的元素的集合。则 \(\mathrm{W}' = \mathrm{W} \otimes_k \overline{k}\)，其中 \(\mathrm{W} = \bigcap_{g \in \mathbb{G}} \mathrm{Ker}(g - 1)\)，并且 \(\mathrm{W}' = \bigcap_{g \in \mathbb{G}} \mathrm{Ker}(g - 1) \otimes 1\)。若 \(\mathrm{V}_1\) 表示 \(\overline{\mathrm{V}}\) 的在 \(\overline{\mathbf{G}}\) 下稳定的非零向量子空间集合中的极小元，则由证明的 (a) 部分，\(\mathrm{V}_1 \subset \mathrm{W}'\)；因此 \(\mathrm{W} \neq \{0\}\)，这就证明了 (i)。
\item
  对 dim V 作归纳，由 (i) 可得，存在 V 的一个在 G 下稳定的向量子空间递增序列 \(\left(\mathrm{V}_{1},\mathrm{V}_{2},\ldots,\mathrm{V}_{n}\right)\)，使得 \(V_{n}=V\)，并且由 G 典范地导出的 \(V_{i}/V_{i-1}\) 的自同构群对所有 i 都退化为 \(\{1\}\)（约定 \(V_{r}=\{0\}\) 对 \(r\leqslant0\) 成立）。这首先推出 (ii)，继而推出 (iii)（第二章第 4 节第 6 号注）。
\end{enumerate}

推论 1. 设 G 是有限维连通实或复李群。G 为幂零的充要条件是 Ad G 的每个元素都是幂零的。

若 Ad G 的每个元素都是幺幂的，则 Ad G 是幂零的（命题 18），因而作为 Ad G 的中心扩张的 G 是幂零的。若 G 是幂零的，则 L(G) 是幂零的，所以对所有 \(x \in \mathrm{L}(G)\)，ad x 都是幂零的，从而 \(\mathrm{Ad}(\exp x) = \exp \mathrm{ad} x\) 是幺幂的；但 G 的每个元素都具有 \(\exp x\) 的形式，其中 \(x \in \mathrm{L}(G)\)（命题 14）。

推论 2. \(\mathbf{GL}(n, \mathbf{K})\) 的每个由幺幂元素组成的解析子群都是单连通李子群。

这由命题 13(ii)、命题 18(ii) 以及严格下三角群是单连通的这一事实推出。

\subsection{可解李群}

命题 19. 设 G 为有限维李群。L(G) 为可解的充要条件是 G 具有一个可解开子群。

证明类似于第 5 号命题 12 的证明。

命题 20. 设 G 是 R 或 C 上维数为 n 的单连通可解李群，且 g = L(G)。设 \((\mathfrak{g}_{n}, \mathfrak{g}_{n-1}, \ldots, \mathfrak{g}_{0})\) 是 g 的子代数序列，其维数依次为 n, n-1, \ldots, 0，并满足 \(g_{i-1}\) 是 \(g_{i}\) 的理想（当 i = n, n-1, \ldots, 1 时）。† 设 \(G_{i}\) 是对应于 \(g_{i}\) 的 G 的积分子群。设 \(x_{i}\) 是 \(g_{i}\) 中不属于 \(g_{i-1}\) 的向量。设 \(\phi_{i}\) 是映射

\[
\left(\lambda_ {1}, \lambda_ {2}, \dots , \lambda_ {i}\right) \mapsto \left(\exp \lambda_ {1} x _ {1}\right) \left(\exp \lambda_ {2} x _ {2}\right) \dots \left(\exp \lambda_ {i} x _ {i}\right)
\]

从 \(\mathbf{K}^i\) 到 G。则 \(\phi_n\) 是解析流形的一个同构，且 \(\phi_i(\mathbf{K}^i) = \mathbf{G}_i\) 对所有 \(i\) 成立。

当 \(n = 0\) 时本命题显然成立。我们对 \(n\) 作归纳。设 H 是 G 中满足 \(\mathrm{L(H)} = \mathrm{K}x_n\) 的积分子群。由第 6 节第 6 号命题 14 的推论 1，H 和 \(G_{n-1}\) 都是 G 的单连通李子群，并且作为李群，G 是 H 与 \(G_{n-1}\) 的半直积。因此 \(\lambda \mapsto \exp (\lambda x_n)\) 是从 K 到 H 的同构，而由归纳假设，映射

\[
\left(\lambda_ {1}, \lambda_ {2}, \dots , \lambda_ {n - 1}\right) \mapsto \left(\exp \lambda_ {1} x _ {1}\right) \left(\exp \lambda_ {2} x _ {2}\right) \dots \left(\exp \lambda_ {n - 1} x _ {n - 1}\right)
\]

是从解析流形 \(\mathbf{K}^{n - 1}\) 到解析流形 \(\mathbf{G}_{n - 1}\) 的同构，并将 \(\mathbf{K}^i\times \{0\}\) 映到 \(\mathbf{G}_i\)，其中 \(i = 1,2,\ldots ,n - 1\)。命题得证。

命题 21. 设 G 是 R 或 C 上的有限维单连通可解李群，M 是 G 的一个积分子群。则 M 是 G 的李子群且是单连通的。

继续使用命题 20 中的记号 \(n\)、\(\mathfrak{g}\)、\(\mathfrak{g}_{i}\)、\(x_{i}\)、\(\phi\)，但对 \(x_{i}\) 附加如下条件：设 \(i_{p} > i_{p-1} > \cdots > i_{1}\) 是这些整数 \(i\)，并且它们满足 \(\mathrm{L}(M) \cap \mathfrak{g}_{i} \neq \mathrm{L}(M) \cap \mathfrak{g}_{i-1}\)；于是取

\[
x _ {i _ {k}} \in \mathrm{L} (\mathrm{M}) \cap \mathfrak {g} _ {i _ {k}}
\]

当 \(k = 1, 2, \ldots, p\) 时。对 \(n\) 作归纳，容易看出 \((x_{i_p}, x_{i_{p-1}}, \ldots, x_{i_1})\) 是 \(\mathrm{L}(\mathbf{M})\) 的一个基。设 \(\mathbf{N}\) 是一个单连通李群，使得存在一个同构 \(h\)，它从 \(\mathrm{L}(\mathbf{N})\) 到 \(\mathrm{L}(\mathbf{M})\)。令

\[
y _ {p} = h ^ {- 1} (x _ {i _ {p}}), \dots , y _ {1} = h ^ {- 1} (x _ {i _ {1}}).
\]

由命题 20，映射

\[
\left(\lambda_ {1}, \lambda_ {2}, \dots , \lambda_ {p}\right) \mapsto \left(\exp \lambda_ {1} y _ {1}\right) \left(\exp \lambda_ {2} y _ {2}\right) \dots \left(\exp \lambda_ {p} y _ {p}\right)
\]

是从流形 \(K^{p}\) 到流形 N 的同构。存在从 N 到 G 的李群态射 \(\tau\)，满足 \(h = L(\tau)\) 且 \(\tau(N) = M\)（\(\S 6\)，第 2 号命题 1 的推论 1）。因此 M 是 G 中形如

\[
\begin{array}{r l} \tau ((\exp \lambda_ {1} y _ {1}) \dots (\exp \lambda_ {p} y _ {p})) & = \exp (\lambda_ {1} L (\tau) y _ {1}) \dots \exp (\lambda_ {p} L (\tau) y _ {p}) \\ & = \exp (\lambda_ {1} x _ {i _ {1}}) \dots \exp (\lambda_ {p} x _ {i _ {p}}). \end{array}
\]

因此 \(\mathbf{M} = \phi (\mathbf{T})\)，其中 \(\mathbf{T}\) 是 \(\mathbf{K}^n\) 的一个向量子空间。

命题 22. 设 K = R 或 C。设 V 是有限维向量空间，G 是 GL(V) 的连通可解子群。假设 G 的恒等表示是单的。

\begin{enumerate}
\def\labelenumi{(\roman{enumi})}
\item
  若 \(\mathbf{K} = \mathbf{R}\)，则 \(\dim V \leqslant 2\)，且 \(G\) 是交换的。
\item
  若 K = C，则 dim V = 1。
\item
  设 \(\mathbf{K} = \mathbf{R}\)。则 G 在 \(\mathbf{GL}(V)\) 中的闭包 H 是 \(\mathbf{GL}(V)\) 的连通可解李子群（第 1 号命题 1 的推论 2）。因此 L(H) 是可解的（命题 19）。L(G) 的恒等表示是单的（第 \(\S 6\)，第 5 号命题 13 的推论 2）。所以 dim V \(\leqslant 2\)，且 L(G) 是交换的（第一章第 \(\S 5\)，定理 1 的推论 1 和 4）。因此 G 是交换的。
\item
  设 K = C。令 W 为 V 的非零实向量子空间中在 G 下稳定的集合的一个极小元。由于 G 的恒等表示是单的，由 W 生成的 V 的复向量子空间等于 V。由 (i)，G\textbar W 是交换的。因此 G 是交换的。于是 G 的每个元素都是一个伸缩变换（《代数》第八章第 4 节命题 2 的推论 I），从而 dim V = 1。
\end{enumerate}

推论。设 V 是维数 \textgreater0 的复向量空间，G 是 GL(V) 的连通可解子群。

\begin{enumerate}
\def\labelenumi{(\roman{enumi})}
\item
  存在一个非零元素 \(v\)，属于 \(V\)，使得 \(gv \in \mathbf{C}v\) 对所有 \(g \in G\) 成立。
\item
  存在一个基 \(\mathbf{B}\)，作为 \(\mathbf{V}\) 的基，使得对所有 \(g \in \mathbf{G}\)，\(g\) 关于 \(\mathbf{B}\) 的矩阵是下三角的。
\end{enumerate}

令 \(V_{1}\) 为 V 的非零向量子空间中在 G 下稳定的集合的一个极小元。由命题 22(ii)，\(\dim V_{1} = 1\)。这就证明了 (i)。再对 \(\dim V\) 作归纳，可得存在 V 的在 G 下稳定的向量子空间递增序列 \((\mathrm{V}_1, \mathrm{V}_2, \ldots, \mathrm{V}_n)\)，使得 \(\dim \mathrm{V}_{i+1} / \mathrm{V}_i = 1\) 对所有 \(i < n\) 成立，且 \(\mathrm{V}_n = \mathrm{V}\)；从而得到 (ii)。

\subsection{李群的根基}

命题 23. 设 G 是有限维实或复李群，r 是 L(G) 的根基（第一章第 5 节定义 2），n 是 L(G) 的最大幂零理想（第一章第 4 节第 4 号）。设 R（分别地，N）是李代数为 r（分别地，n）的 G 的积分子群。则 R（分别地，N）是 G 的可解（分别地，幂零）李子群，并且在 G 的每个连续自同构下都不变。G 的每个连通可解（分别地，幂零）正规子群都包含于 R（分别地，N）中。

群 R 是可解的（第 6 节命题 19）。设 K = R。令 G' 为 G 的连通可解正规子群。则 \(\overline{G'}\) 是 G 的连通可解（第 1 号命题 1 的推论 2）正规李子群（第 8 节第 2 号定理 2）。因此 L \((\overline{G'})\) 是 L(G) 的可解理想，从而 L \((\overline{G'}) \subset r\)，且 \(\overline{G'} \subset R\)。特别地，\(\overline{R} \subset R\)，所以 R 是闭的，因而是 G 的李子群。设 K = C。令 H 为 G 的底层实李子群。若 \(r'\) 是 L(H) 的根基，则 \(i'r\) 是 L(H) 的可解理想，从而 \(r' = i'r'\)；于是 \(r \subset r' \subset r\)，由上面所述，R 在 H 中闭，因而在 G 中闭；所以 R 是 G 的李子群。G 的每个连通可解正规子群都是 H 的连通可解正规子群，因而包含于 R 中。因此我们已经证明，当 K = C 或 K = R 时，R 是 G 的最大连通可解正规子群；所以 R 在 G 的每个连续自同构下都不变。关于 N 的证明完全类似。

定义 1. 设 G 是有限维实或复李群。G 的根基是 G 的最大连通可解正规子群。

注。即使 G 是连通的，G 也可能有不包含于 G 的根基中的可解正规子群。

命题 24. 设 K = R 或 C。令 \(G_{1}\)、\(G_{2}\) 为两个有限维连通李群，\(R_{1}\) 和 \(R_{2}\) 为它们的根基，\(\phi\) 为从 \(G_{1}\) 到 \(G_{2}\) 的满射态射。则 \(\phi(R_{1}) = R_{2}\)。

由第 3 节第 8 号命题 28，\(\mathrm{L}(\phi)\) 是满射。因此 \(\mathrm{L}(\phi)(\mathrm{L}(\mathbb{R}_1)) = \mathrm{L}(\mathbb{R}_2)\)（第一章第 6 节命题 2 的推论 3）。令 \(i\) 为从 \(\mathbb{R}_1\) 到 \(\mathbb{G}_1\) 的典范嵌入。则 \(\phi \circ i\) 的像是 \(\mathbb{R}_2\)（第 6 节第 2 号命题 1 的推论 I）。

命题 25. 设 \(\mathbf{K} = \mathbf{R}\) 或 \(\mathbf{C}\)。令 \(\mathbf{G}_1, \mathbf{G}_2\) 为有限维连通李群，\(\mathbf{R}_1\) 和 \(\mathbf{R}_2\) 为它们的根基。\(\mathbf{G}_1 \times \mathbf{G}_2\) 的根基是 \(\mathbf{R}_1 \times \mathbf{R}_2\)。这由第一章第 5 节命题 4 推出。

\subsection{半单李群}

命题 26. 设 \(G\) 是有限维连通实或复李群。以下条件等价：

\begin{enumerate}
\def\labelenumi{(\roman{enumi})}
\item
  L(G) 是半单的；
\item
  G 的根基是 \{e\}；
\item
  \(\mathbf{G}\) 的每个正规交换积分子群都等于 \(\{e\}\)。
\end{enumerate}

条件 (ii) 意味着 L(G) 的根基是 \(\{0\}\)，所以 (i) \(\Leftrightarrow\) (ii)（第一章第 6 节定理 1）。(i) 与 (iii) 的等价性由第 6 节第 6 号命题 14 得出。

定义 2. 若一个连通实或复李群是有限维的且满足命题 26 的条件，则称其为半单李群。

注 1. 设 G 是有限维连通实或复李群。若 G 不是半单的，则 G 有一个交换连通李子群 \(G'\)，它在每个连续自同构下都不变，且 \(G' \neq \{\varepsilon\}\)。事实上，令 n 为 L(G) 的最大幂零理想；则 \(n \neq \{0\}\)，相应的解析子群 N 是一个李子群，并在 G 的每个连续自同构下都不变（第 7 号命题 23）；N 的中心 G' 具有所需性质。

命题 27. 设 \(G\) 是有限维连通实或复李群。以下条件等价：

\begin{enumerate}
\def\labelenumi{(\roman{enumi})}
\item
  L(G) 是单的；
\item
  \(\mathbf{G}\) 的正规积分子群只有 \(\{e\}\) 和 \(\mathbf{G}\)，并且 \(\mathbf{G}\) 不是交换的。
\end{enumerate}

这由第 6 节第 6 号命题 14 得出。

定义 3. 若一个连通实或复李群是有限维的且满足命题 27 的条件，则称其为几乎单李群。

命题 28. 设 \(\mathbf{G}\) 是单连通实或复李群。以下条件等价：

\begin{enumerate}
\def\labelenumi{(\roman{enumi})}
\item
  G 是半单的；
\item
  G 同构于有限个几乎单群的乘积。
\end{enumerate}

若 G 是有限个几乎单李群的乘积，则 \(\mathrm{L}(\mathrm{G})\) 是有限个单李代数的乘积，因而是半单的。若 G 是半单的，则 \(\mathrm{L}(\mathrm{G})\) 同构于单李代数 \(L_{1},\ldots,L_{n}\) 的乘积。令 \(G_{i}\) 是李代数为 \(L_{i}\) 的单连通李群，因而它是几乎单的。于是 G 与 \(G_{1}\times\cdots\times G_{n}\) 都是单连通的，且具有同构的李代数，所以二者同构。

引理 1. 设 \(\mathbf{G}\) 是连通拓扑群，\(\mathbf{Z}\) 是其中心，\(\mathbf{Z}'\) 是 \(\mathbf{Z}\) 的离散子群。则 \(\mathbf{G}/\mathbf{Z}'\) 的中心是 \(\mathbf{Z}/\mathbf{Z}'\)。

设 \(y\) 是 G 的一个元素，其模 \(\mathbf{Z}'\) 的类是

G/Z' 的中心元素。令 \(\phi\) 为从 G 到 G 的映射 \(g \mapsto gyg^{-1}y^{-1}\)。则 \(\phi(G)\) 是连通的且包含于 Z' 中，所以 \(\phi(G) = \phi(\{e\}) = \{e\}\)。因此 \(y \in Z\)。

命题 29. 设 G 是半单连通实或复李群。

\begin{enumerate}
\def\labelenumi{(\roman{enumi})}
\item
  G = (G, G)。
\item
  G 的中心 Z 是离散的。
\item
  G/Z 的中心是 \{e\}。
\end{enumerate}

断言 (i) 由命题 4 的第 2 号推论和第一章第 6 节定理 1 得出。

断言 (ii) 由第 6 节第 4 号命题 10 的推论 4 以及第一章第 6 节第 1 号注 2 得出。

断言 (iii) 由 (ii) 和引理 1 得出。

命题 30. (i) 设 g 是半单实或复李代数。则 Int g 是 Aut g 的单位元分支。

\begin{enumerate}
\def\labelenumi{(\roman{enumi})}
\setcounter{enumi}{1}
\tightlist
\item
  设 \(\mathbf{G}\) 是半单连通实或复李群。\(\mathbf{G}\) 的伴随群是 \(\operatorname{Aut} \mathrm{L}(\mathrm{G})\) 的单位元分支。其中心退化为单位元。
\end{enumerate}

g 的每个导子都是内导子（第一章第 6 节命题 1 的推论 3），所以 L(Int g) = L(Aut g)，这就证明了 (i)。(ii) 的第一个断言由 (i) 得出，第二个断言由命题 29(iii) 以及第 6 节第 4 号命题 10 的推论 4(ii) 得出。

注 2. 设 g 是复半单李代数，\(g_{0}\) 是其底层实李代数。则 \(\operatorname{Aut}(g)\) 在 \(\operatorname{Aut}(g_{0})\) 中是开的，因为 \(\operatorname{Int}(g_{0}) \subset \operatorname{Aut}(g)\)。

命题 31. 设 G 是有限维单连通实或复李群，R 是其根基。存在 G 的一个半单单连通李子群 S，使得作为李群，G 是 S 与 R 的半直积。若 S' 是 G 的一个半单积分子群，则存在 L(G) 的幂零根基中的 x，使得

\[
(\operatorname{Ad} \exp x) (\mathrm{S} ^ {\prime}) \subset \mathrm{S}.
\]

这由第 6 节第 6 号命题 14 的推论 1 以及第一章第 6 节定理 5 和推论 1 得出。

引理 2. 设 G 是一个群（分别地，拓扑群），G' 是 G 的正规子群，V 是交换域 k（分别地，K）上的有限维向量空间，\(\rho\) 是 G 在 V 上的线性表示（分别地，连续线性表示），且 \(\rho' = \rho|G'\)。

\begin{enumerate}
\def\labelenumi{(\roman{enumi})}
\item
  若 \(\rho\) 是半单的，则 \(\rho'\) 是半单的。
\item
  若 \(\rho'\) 是半单的，且每个有限维线性 \(k\)-表示（分别地，连续线性 K-表示）都是 \(\mathrm{G} / \mathrm{G}'\)（分别地，\(\mathrm{G} / \overline{\mathrm{G}}'\)）的半单表示，则 \(\rho\) 是半单的。
\end{enumerate}

设 \(\rho\) 是半单的；我们证明 \(\rho'\) 是半单的。只需考虑 \(\rho\) 是单的情形。令 \(V'\) 为一个极小非零 \(G'\)-子模，它是 \(V\) 的子模。对所有 \(g \in G\)，\(\rho(G') \rho(g)V' = \rho(g)\rho(G')V' = \rho(g)V'\)，换言之，\(\rho(g)\mathrm{V}'\) 在 \(\rho(\mathrm{G}')\) 下稳定；若 \(\mathrm{V}''\) 是 \(\mathrm{G}'\) 的一个子模且包含于 \(\rho(g)\mathrm{V}'\)，则 \(\rho(g)^{-1}\mathrm{V}''\) 是 \(\mathrm{G}'\) 的一个子模且包含于 \(\mathrm{V}'\)，因而 \(\mathrm{V}''\) 等于 \(\{0\}\) 或 \(\rho(g)\mathrm{V}'\)。所以，对所有 \(g \in \mathrm{G}\)，\(\rho(g)\mathrm{V}'\) 都是一个单的 \(\mathrm{G}'\)-模。但 \(\sum_{g \in \mathrm{G}} \rho(g)\mathrm{V}'\) 是 \(\mathrm{V}\) 的非零子-G-模，从而 \(\mathrm{V} = \sum_{g \in \mathrm{G}} \rho(g)\mathrm{V}'\)。因此 \(\rho'\) 是半单的。

设 \(\rho'\) 是半单的。令 \(W\) 是 \(V\) 的一个非零子-G-模。由于 \(\rho'\) 是半单的，存在一个 \(f_0\)，它是从 \(V\) 到 \(W\) 的投影算子，并且与 \(\rho'(G)\) 对易。令 \(E\) 为满足以下条件的集合：其中的 \(f \in \mathcal{L}(V, V)\) 与 \(\rho(G')\) 对易，将 \(V\) 映入 \(W\)，且其在 \(W\) 上的限制是伸缩变换；对 \(f \in E\)，令 \(\alpha(f)\) 表示伸缩变换 \(f|W\) 的比值。于是 \(f_0 \in E\)，且 \(\alpha(f_0) = 1\)。显然，\(\alpha\) 是 \(E\) 上的线性形式。令 \(F = \text{Ker } \alpha\)，它是 \(E\) 的一个超平面。对 \(f \in E\) 和 \(g \in G\)，记 \(\sigma(g)f = \rho(g) \circ f \circ \rho(g)^{-1}\)；则 \(\sigma(g)f\) 将 \(V\) 映入 \(W\)，且其在 \(W\) 上的限制是比值为 \(\alpha(f)\) 的伸缩变换；若 \(g' \in G'\)，则

\[
\begin{array}{r l} \sigma (g) f \circ \rho (g ^ {\prime}) & = \rho (g) \circ f \circ \rho (g) ^ {- 1} \circ \rho (g ^ {\prime}) \\ & = \rho (g) \circ f \circ \rho (g ^ {- 1} g ^ {\prime} g) \circ \rho (g ^ {- 1}) \\ & = \rho (g) \circ \rho (g ^ {- 1} g ^ {\prime} g) \circ f \circ \rho (g ^ {- 1}) \\ & = \rho (g ^ {\prime}) \circ \rho (g) \circ f \circ \rho (g ^ {- 1}) \\ & = \rho (g ^ {\prime}) \circ \sigma (g) f. \end{array}
\]

因此 \(\sigma(g)f\in\mathbf{E}\)。于是 \(\sigma\) 是 G 在 E 上的线性 \(k\)-表示（分别地，连续线性 K-表示），并使 F 稳定。于是 \(\sigma(g)=\mathrm{Id}_{\mathbf{E}}\) 当 \(g\in\mathbf{G}'\) 时成立，在拓扑情形下进而当 \(g\in\overline{\mathbf{G}'}\) 时也成立。假设每个有限维 \(k\)-线性表示（分别地，连续 K-线性表示）作用在 G/G'（分别地，G/ \(\overline{\mathbf{G}'}\)）上都是半单的。于是 E 中存在 F 的一个在 G 下稳定的补空间。换言之，存在 \(f\in\mathbf{E}\)，使得 \(\alpha(f)=1\)，且 \(f\) 在 G 下不变。于是它是从 V 到 W 的投影算子，并且对 \(g\in\mathbf{G}\)，\(\rho(g)\circ f\circ\rho(g^{-1})=f\)，即 \(f\) 与 \(\rho(\mathbf{G})\) 对易。因此 \(\rho\) 是半单的。

定理 1. 设 G 是有限维实或复李群，\(G_{0}\) 是其单位元分支，R 是其根基，r 是 L(G) 的根基；假设 \(G/G_{0}\) 是有限的。设 \(\rho\) 是 G 的有限维解析线性表示。以下条件等价：

\begin{enumerate}
\def\labelenumi{(\roman{enumi})}
\item
  \(\rho\) 是半单的；
\item
  \(\rho|G_{0}\) 是半单的；
\item
  \(\rho|R\) 是半单的；
\item
  L(ρ) 是半单的；
\item
  L(ρ)\textbar r 是半单的。
\end{enumerate}

(i) \(\Leftrightarrow\) (ii) 由引理 2 和《积分》第七章第 3 节命题 1 得出。(ii) \(\Leftrightarrow\) (iv) 以及 (iii) \(\Leftrightarrow\) (v) 由第 6 节第 5 号命题 13 的推论 2 得出。(iv) \(\Leftrightarrow\) (v) 由第一章第 6 节定理 4 得出。

推论 1. 设 \(\rho, \rho_{1}, \rho_{2}\) 是 G 的有限维半单解析线性表示，n 是一个整数 \(\geqslant 0\)。则 \(\rho_{1} \otimes \rho_{2}\)、\(T^{n}\rho\)、\(S^{n}\rho\)、\(\bigwedge^{n}\rho\)（附录）都是半单的。

\(\rho_{1} \otimes \rho_{2}\) 的半单性由定理 1 和第一章第 6 节定理 4 的推论 1 得出。\(T^{n}_{\rho}, S^{n}_{\rho}, \bigwedge^{n}_{\rho}\) 的半单性由 \(\rho_{1} \otimes \rho_{2}\) 的半单性得出。

我们将在稍后看到，若 \(k\) 是特征为 0 的交换域，\(\Gamma\) 是一个群，且 \(\rho_1\) 和 \(\rho_2\) 是有限维半单线性 \(k\)-表示，其群为 \(\Gamma\)，则 \(\rho_1 \otimes \rho_2\) 是半单的。

推论 2. 设 \(\rho\) 是 \(G\) 在向量空间 \(V\) 上的有限维半单解析线性表示，\(S\) 是 \(V\) 的对称代数，\(S^{\mathfrak{G}}\) 是 \(S\) 的由在 \((\mathsf{S}(\rho))(G)\) 下不变的元素组成的子代数。则 \(S^{\mathfrak{G}}\) 是有限生成代数。

这由定理 1、第一章第 6 节定理 6(a) 以及《交换代数》第五章第 1 节定理 2 得出。

推论 3. 设 G 是实或复李群，\(G_{0}\) 是其单位元分支。假设 \(G_{0}\) 是半单的且 \(G/G_{0}\) 是有限的。则 G 的每个有限维解析线性表示都是半单的。

命题 32. 设 G 是有限维连通实李群。假设 L(G) 是约化的。以下条件等价：

\begin{enumerate}
\def\labelenumi{(\roman{enumi})}
\item
  \(\mathbf{G} / \overline{\mathbf{D}^1}\mathbf{G}\) 是紧的；
\item
  （分别地，(ii')）\(G\) 在复（分别地，实）向量空间上的每个有限维解析线性表示都是半单的。
\item
  \(\Rightarrow\) (ii')：假设 \(G / \overline{D}^1 G\) 是紧的。则 \(G / \overline{D}^1 G\) 在有限维实向量空间上的每个连续线性表示都是半单的（《积分》第七章第 3 节命题 1）。设 \(\rho\) 是 \(G\) 在实向量空间上的有限维解析线性表示。则 \(\rho |D^1 G\) 是解析的，\(D^1 G\) 是半单的（第一章第 6 节命题 5），所以 \(\rho |D^1 G\) 是半单的（定理 1 的推论 3）。因此 \(\rho\) 是半单的（引理 2）。
\end{enumerate}

同理可见 (i) \(\Rightarrow\) (ii)。

\begin{enumerate}
\def\labelenumi{(\roman{enumi})}
\tightlist
\item
  \(\Rightarrow\) (ii)：假设 \(G / \overline{D^1} G\) 不是紧的，因而同构于形如 \(\mathbf{R}^p\times \mathbf{T}^q\) 的群，其中 \(p > 0\)（第 6 节第 4 号命题 11(ii)）。于是存在从 \(G / \overline{D^1} G\) 到 \(\mathbf{R}\) 的满射态射，从而存在从 \(G\) 到 \(\mathbf{R}\) 的满射态射。映射
\end{enumerate}

\[
g \mapsto \sigma (g) = \left( \begin{array}{c c} 1 & 0 \\ \rho (g) & 1 \end{array} \right)
\]

是 G 在 \(R^{2}\) 上的解析线性表示，但不是半单的，因为在 \(\mathbf{R}^2\) 的一维向量子空间中，唯一在 \(\sigma(G)\) 下稳定的是 \(\mathbf{R}(0,1)\)。

同理可见 (ii) \(\Rightarrow\) (i)。

命题 33. 设 \(G\) 是连通分支数有限的有限维复李群，\(\rho\) 是 \(G\) 的有限维解析线性表示，\(G'\) 是实李群 \(G\) 的一个积分子群，且 \(L(G')\) 在 \(L(G)\) 上作为 \(C\)-向量空间生成。则 \(\rho\) 为半单的充要条件是 \(\rho |G'\) 为半单的。

令 \(\rho' = \rho|G'\)。\(\rho\)（分别地，\(\rho'\)）为半单的充要条件是 \(L(\rho)\)（分别地，\(L(\rho')\)）为半单的（定理 1）。令 \(V\) 为 \(\rho\) 的空间。\(V\) 的一个向量子空间在 \(L(\rho)(L(G))\) 下稳定的充要条件是它在 \(L(\rho')(L(G'))\) 下稳定。因此命题得证。

\section{李群的自同构群}

本段中假设 K 的特征为零。

\subsection{无穷小自同构}

引理 1. 设 G 是李群，\(\alpha\) 是 G 上的向量场。对所有 \(g \in \mathbf{G}\)，令

\[
\beta (g) = \alpha (g) g ^ {- 1} \in \mathrm{L} (\mathrm{G}).
\]

以下条件等价：

\begin{enumerate}
\def\labelenumi{(\roman{enumi})}
\item
  \(\alpha\) 是从群 G 到群 T(G) 的同态；
\item
  对所有 \(g, g'\) 属于 \(\mathbf{G}\)，都有 \(\alpha(gg') = \alpha(g)g' + g\alpha(g')\)；
\item
  对所有 \(g, g'\) 属于 \(\mathbf{G}\)，都有 \(\beta(gg') = \beta(g) + (\mathrm{Ad}g)\beta(g')\)。
\end{enumerate}

条件 (i) 意味着，对所有 \(g, g'\) 属于 \(G\)，在群 \(T(G)\) 中有：

\[
\beta (g) g \beta (g ^ {\prime}) g ^ {\prime} = \beta (g g ^ {\prime}) g g ^ {\prime}
\]

或者

\[
\beta (g) ((\mathrm{Ad} g) \beta (g ^ {\prime})) g g ^ {\prime} = \beta (g g ^ {\prime}) g g ^ {\prime}.
\]

但是 \(\beta(g)\) 与 (Ad \(g)\beta(g')\) 在 T(G) 中的乘积，正是 \(\beta(g)\) 与 (Ad \(g)\beta(g')\) 在 L(G) 中的和（第 2 节第 1 号命题 2）。因此 (i) \(\Leftrightarrow\) (iii)。另一方面，条件 (ii) 可写成 \(\beta(gg')gg' = \beta(g)gg' + g\beta(g')g'\)，或者

\[
\beta (g g ^ {\prime}) = \beta (g) + (\mathrm{Ad} g) \beta (g ^ {\prime})
\]

从而 (ii) \(\Leftrightarrow\) (iii)。

定义 1. 设 G 是李群。G 的无穷小自同构是 G 上满足引理 1 条件的任一解析向量场。

引理 2. 设 \(\mathbf{K}'\) 是 \(\mathbf{K}\) 的非离散闭子域，A 是 \(\mathbf{K}'\)-流形，B 和 C 是 K-流形，且 \(f\) 是 \(\mathbf{K}'\)-解析映射，从 \(\mathbf{A} \times \mathbf{B}\) 到 C。假设对所有 \(a \in \mathbf{A}\)，映射 \(b \mapsto f(a, b)\) 从 \(\mathbf{B}\) 到 \(\mathbf{C}\) 都是 K-解析的。于是对所有 \(t \in \mathbf{TA}\)，从 TB 到 TC 的映射 \(u \mapsto (\mathrm{Tf})(t, u)\) 是 K-解析的。

固定 \(t \in \mathrm{TA}\)，记 \(g(u) = (\mathrm{Tf})(t,u)\)。显然，\(g\) 是 K'-解析的。由《可微与解析流形》R，5.14.6，只需证明 \(g\) 的切映射是 K-线性的。可以假定 A、B、C 分别是 K'、K、K 上完备可赋范空间 E、F、G 中 0 的开邻域，且 \(t\) 是 A 在 0 处的切向量。将 TA、TB、TC 分别与 A × E、B × F、C × G 识别，并将 \(t\) 与 E 的一个元素识别。于是对所有 \((x,y) \in \mathrm{TB} = \mathrm{B} \times \mathrm{F}\)。

\[
g (x, y) = \left(f (0, x), \left(\mathrm{D} _ {1} f\right) (0, x) (t) + \left(\mathrm{D} _ {2} f\right) (0, x) (y)\right).
\]

将 \(\mathrm{T(B\times F)}\) 与 \((\mathbf{B}\times \mathbf{F})\times (\mathbf{F}\times \mathbf{F})\) 识别，将 \(\mathrm{T(C\times G)}\) 与 \((\mathbf{C}\times \mathbf{G})\times (\mathbf{G}\times \mathbf{G})\) 识别。于是，对所有

\[
((x, y), (h, k)) \in \mathrm{T} (\mathrm{B} \times \mathrm{F}) = (\mathrm{B} \times \mathrm{F}) \times (\mathrm{F} \times \mathrm{F}),
\]

\((\mathrm{Tg})((x,y),(h,k)) = ((a,b),(c,d)),\) 其中

\[
\begin{array}{r l} & {a = f (0, x),} \\ & {b = (\mathrm{D} _ {1} f) (0, x) (t) + (\mathrm{D} _ {2} f) (0, x) (y), \qquad c = (\mathrm{D} _ {2} f) (0, x) (h),} \\ & {d = (\mathrm{D} _ {2} \mathrm{D} _ {1} f) (0, x) (t, h) + (\mathrm{D} _ {2} \mathrm{D} _ {2} f) (0, x) (y, h) + (\mathrm{D} _ {2} f) (0, x) (k).} \end{array}
\]

现在固定 \((x,y)\in \mathbf{B}\times \mathbf{F}\)。我们需要证明映射 \((h,k)\mapsto (c,d)\) 从 \(\mathbf{F}\times \mathbf{F}\) 到 \(\mathbf{G}\times \mathbf{G}\) 是 K-线性的。由于从 B 到 C 的映射 \(x\mapsto f(0,x)\) 是 K-解析的，映射

\[
\begin{array}{c} (h, k) \mapsto (\mathrm{D} _ {2} f) (0, x) (h), \qquad (h, k) \mapsto (\mathrm{D} _ {2} \mathrm{D} _ {2} f) (0, x) (y, h), \\ (h, k) \mapsto (\mathrm{D} _ {2} f) (0, x) (k) \end{array}
\]

是 K-线性的。另一方面

\[
\left(\mathrm{D} _ {2} \mathrm{D} _ {1} f\right) (0, x) (t, h) = \lim _ {\lambda \in K ^ {*}, \lambda \rightarrow 0} \lambda^ {- 1} \left(\left(\mathrm{D} _ {2} f\right) (\lambda t, x) (h) - \left(\mathrm{D} _ {2} f\right) (0, x) (h)\right)
\]

并且，对固定的 \(\lambda\)，映射 \(x \mapsto f(\lambda t, x)\) 是 K-解析的，所以映射 \(h \mapsto (\mathrm{D}_2f)(\lambda t, x)(h)\) 是 K-线性的。

命题 1. 设 \(\mathbf{K}'\) 是 \(\mathbf{K}\) 的非离散闭子域，\(\mathbf{G}\) 是 \(\mathbf{K}\) 上的李群，\(\mathbf{V}\) 是 \(\mathbf{K}'\) 上的流形，且 \((v, g) \mapsto vg\) 是 \(\mathbf{K}'\)-解析映射，从 \(\mathbf{V} \times \mathbf{G}\) 到 \(\mathbf{G}\)。假设对所有 \(v \in \mathbf{V}\)，映射 \(g \mapsto vg\) 从 \(\mathbf{G}\) 到 \(\mathbf{G}\) 都是 \(\mathbf{G}\) 的自同构。设 \(\varepsilon\) 是 \(\mathbf{V}\) 的一个元素，满足 \(\varepsilon g = g\) 对所有 \(g \in \mathbf{G}\) 成立，且 \(a \in T_{\varepsilon}(\mathbf{V})\)。则 \(g \mapsto ag\) 是 \(\mathbf{G}\) 上的向量场，也是 \(\mathbf{G}\) 的无穷小自同构。

对 \(v \in \mathrm{V}\)、\(g_1 \in \mathrm{G}\)、\(g_2 \in \mathrm{G}\)，有 \(v(g_1g_2) = (vg_1)(vg_2)\)。因此，对 \(u_1 \in \mathrm{TG}\)、\(u_2 \in \mathrm{TG}\)，有 \(a(u_1u_2) = (au_1)(au_2)\)（第 2 节第 1 号命题 3）。特别地，映射 \(g \mapsto ag\) 从 \(G\) 到 \(TG\) 是群同态。另一方面，由引理 2，这个映射是解析的。

命题 2. 设 G 是实或复李群，\(\alpha\) 是 G 的无穷小自同构。存在 K 在 G 上的解析运算律 \((\lambda, g) \mapsto \phi_{\lambda}(g)\)，具有以下性质：

\begin{enumerate}
\def\labelenumi{(\arabic{enumi})}
\item
  若 D 是相应的无穷小运算律，则 D(1) = α；
\item
  对所有 \(\lambda \in \mathbf{K}\)，\(\phi_{\lambda} \in \operatorname{Aut} G\)。
\end{enumerate}

\begin{enumerate}
\def\labelenumi{(\alph{enumi})}
\item
  对所有 \(\mu > 0\)，令 \(\mathbf{K}_{\mu}\) 为半径为 \(\mu\)、位于 \(\mathbf{K}\) 中的开球。对所有 \(g \in G\)，令 \(\mathcal{F}_g\) 为由 \(f\)（\(\alpha\) 的、在 \(\mathbf{K}_{\mu}\) 上定义的解析积分曲线）组成的集合，且满足 \(f(0) = g\)。由《可微与解析流形》R，9.1.3 和 9.1.5，\(\mathcal{F}_g\) 非空，且 \(\mathcal{F}_g\) 的两个元素在其定义域的交集上重合；令 \(\mu(g)\) 为满足以下条件的 \(\mu\)：\(\mathcal{F}_g\) 中存在定义在 \(\mathbf{K}_{\mu}\) 上的元素；\(\mathcal{F}_g\) 中存在唯一一个定义在 \(\mathbf{K}_{\mu(0)}\) 上的元素；记其为 \(f_g\)。
\item
  设 \(g_1, g_2\) 属于 \(\mathbf{G}\)，\(f_1 \in \mathcal{F}_{g_1}\)，\(f_2 \in \mathcal{F}_{g_2}\)，且 \(f_1\) 与 \(f_2\) 定义在同一个球 \(\mathbf{K}_\mu\) 上。则 \(f_1f_2: \mathbf{K}_\mu \to \mathbf{G}\) 是解析的，且 \((f_1f_2)(0) = g_1g_2\)。另一方面，对所有 \(\lambda \in \mathbf{K}_\mu\)，
\end{enumerate}

\[
\begin{array}{r l} (\mathrm{T} _ {\lambda} (f _ {1} f _ {2})) 1 & = (\mathrm{T} _ {\lambda} f _ {1}) 1 \cdot f _ {2} (\lambda) + f _ {1} (\lambda) \cdot (\mathrm{T} _ {\lambda} f _ {2}) 1 \quad (\S 2, \text { Proposition   7 }) \\ & = \alpha (f _ {1} (\lambda)) f _ {2} (\lambda) + f _ {1} (\lambda) \alpha (f _ {2} (\lambda)) \\ & = \alpha ((f _ {1} f _ {2}) (\lambda)) \quad \text {(Lemma   1)} \end{array}
\]

从而 \(f_{1}f_{2}\in \mathcal{F}_{g_{1}g_{2}}\)。这就证明了 \(\mu (g_1g_2)\geqslant \inf (\mu (g_1),\mu (g_2))\)。

\begin{enumerate}
\def\labelenumi{(\alph{enumi})}
\setcounter{enumi}{2}
\tightlist
\item
  由《可微与解析流形》R，9.1.4 和 9.1.5，存在 G 中 \(e\) 的一个邻域 V，使得 \(\sigma = \inf_{g \in \mathbf{V}} \mu(g) > 0\)。令 \(h \in \mathbf{G}\)，C 为其连通分支。由 (b)，对所有 \(h' \in \mathbf{C}\)，\(\mu(h') \geqslant \inf(\sigma, \mu(h)) > 0\)。另一方面，函数 \(f_h'\)（其中 \(h' \in \mathbf{C}\)）的值都在 C 中。由《可微与解析流形》R，9.1.4 和 9.1.5，C 中 \(\mu = +\infty\)，最终 G 中 \(\mu = +\infty\)。现在令 \(f_g(\lambda) = \phi_{\lambda}(g)\)，其中 \(g \in \mathbf{G}\) 和 \(\lambda \in \mathbf{K}\) 任意。由《可微与解析流形》R，9.1.4 和 9.1.5，映射 \((\lambda, g) \mapsto \phi_{\lambda}(g)\) 是 K 在 G 上的解析运算律。显然，若 D 是相应的无穷小运算律，则 D(1) = \(\alpha\)。由 (b)，
\end{enumerate}

\[
\phi_ {\lambda} (g _ {1} g _ {2}) = \phi_ {\lambda} (g _ {1}) \phi_ {\lambda} (g _ {2})
\]

对所有 \(\lambda \in \mathbf{K}\)、\(g_{1} \in \mathbf{G}\)、\(g_{2} \in \mathbf{G}\)，均成立。

命题 3. 假设 K 是超度量的。设 G 是紧李群，\(\alpha\) 是 G 的无穷小自同构。存在 K 的一个开子群 I 以及 I 在 G 上的解析运算律 \((\lambda, g) \mapsto \phi_{\lambda}(g)\)，具有以下性质：

\begin{enumerate}
\def\labelenumi{(\arabic{enumi})}
\item
  若 D 是相应的无穷小运算律，则 D(1) = α；
\item
  对所有 \(\lambda \in \mathbf{I}\)，\(\phi_{\lambda} \in \operatorname{Aut} G\)。
\end{enumerate}

由于 G 是紧的，存在 K 的一个开子群 I' 以及 I' 在 G 上的解析运算律 \((\lambda, g) \mapsto \phi_{\lambda}(g)\)，具有本命题的性质 (1)（第 4 节第 7 号定理 6 的推论 2）。记 \(\phi_{\lambda}(g) = f_{g}(\lambda)\)，其中 \(\lambda \in I'\) 和 \(g \in G\)。于是，对 \(g_{1}, g_{2}\) 属于 \(G\) 和 \(\lambda \in I'\)，

\[
\begin{array}{r l} (\mathrm{T} _ {\lambda} (f _ {g _ {1}} f _ {g _ {2}})) 1 & = (\mathrm{T} _ {\lambda} f _ {g _ {1}}) 1 \cdot f _ {g _ {2}} (\lambda) + f _ {g _ {1}} (\lambda) \cdot (\mathrm{T} _ {\lambda} f _ {g _ {2}}) 1 \\ & = \alpha (f _ {g _ {1}} (\lambda)) f _ {g _ {2}} (\lambda) + f _ {g _ {1}} (\lambda) \alpha (f _ {g _ {2}} (\lambda)) \\ & = \alpha (f _ {g _ {1}} (\lambda) f _ {g _ {2}} (\lambda)) \end{array}
\]

并且 \((f_{g_1}f_{g_2})(0) = g_1g_2 = f_{g_1g_2}(0)\)。因此 \(f_{g_1'g_2'}(\lambda) = f_{g_1'}(\lambda)f_{g_2'}(\lambda)\) 对

\[
(g _ {1} ^ {\prime}, g _ {2} ^ {\prime}, \lambda)
\]

在 \((g_1, g_2, 0)\) 的一个邻域内成立（《可微与解析流形》R，9.1.8）。由于 G 是紧的，存在 I' 的一个开子群 I，使得 \(f_{g_1 g_2}(\lambda) = f_{g_1}(\lambda) f_{g_2}(\lambda)\) 对所有 \(g_1 \in G\)、\(g_2 \in G\)、\(\lambda \in I\) 成立。换言之，\(\phi_\lambda \in \operatorname{Aut} G\)（当 \(\lambda \in I\) 时）。

引理 3. 设 G 和 G' 是李群，\(\phi\) 是从 G 到 \(\operatorname{Aut}(\mathbf{G}')\) 的同态。令 \(f(g, g') = (\phi(g))(g')\)，其中 \(g \in \mathbf{G}\)、\(g' \in \mathbf{G}'\)。考虑以下条件：

\begin{enumerate}
\def\labelenumi{(\roman{enumi})}
\item
  \(f\) 是解析的；
\item
  \(f\) 在 \((e_{\mathbf{G}}, e_{\mathbf{G}^{\prime}})\) 的一个邻域内是解析的；
\item
  对所有 \(g' \in G'\)，映射 \(g \mapsto f(g, g')\) 都是解析的。
\end{enumerate}

则 (i) \(\Leftrightarrow\) ((ii) 且 (iii))。若 G 是连通的，则 (i) \(\Leftrightarrow\) (ii)。

显然 (i) 蕴含 (ii) 和 (iii)。设 \(g_0 \in G, g_0' \in G'\)。对所有 \(g \in G, g' \in G'\)，

\[
f \left(g g _ {0}, g ^ {\prime} g _ {0} ^ {\prime}\right) = (\phi (g) \phi (g _ {0})) \left(g ^ {\prime} g _ {0} ^ {\prime}\right) = \phi (g) \left(\phi (g _ {0}) g ^ {\prime}\right)\cdot \phi (g) \left(\phi (g _ {0}) g _ {0} ^ {\prime}\right).
\]

这证明了蕴含 ((ii) 且 (iii)) \(\Rightarrow\) (i)。最后，若 \(\mathbf{G}'\) 是连通的，则 \(\mathbf{G}'\) 由 \(e_{\mathbf{G}'}\) 的每个邻域生成，从而 (ii) \(\Rightarrow\) (iii)。

\subsection{李群的自同构群（实或复情形）}

本号中假设 K = R 或 C。

引理 4. 设 \(\mathbf{H}\) 是有限维单连通李群。

\begin{enumerate}
\def\labelenumi{(\roman{enumi})}
\item
  对所有 \(u \in \operatorname{Aut} \mathbf{L}(\mathbf{H})\)，令 \(\theta(u)\) 为 \(\mathbf{H}\) 的唯一自同构，使得 \(\mathbf{L}(\theta(u)) = u\)。则映射 \((u, g) \mapsto \theta(u)g\) 从 \((\operatorname{Aut} \mathbf{L}(\mathbf{H})) \times \mathbf{H}\) 到 \(\mathbf{H}\) 是解析的。
\item
  设 \(\mathbf{N}\) 是 \(\mathbf{H}\) 的李子群，\(\operatorname{Aut}(\mathbf{H},\mathbf{N})\) 是由 \(v\in \operatorname{Aut}\mathbf{H}\) 组成的集合，其中 \(v(\mathbf{N}) = \mathbf{N}\)。则 \(\theta^{-1}(\operatorname {Aut}(\mathbf{H},\mathbf{N}))\) 是 \(\operatorname {Aut}\mathbf{L}(\mathbf{H})\) 的李子群。
\item
  假设 N 是离散正规子群，从而 G = H/N 的李代数可与 L(H) 识别。对所有 \(w \in Aut G\)，令 \(\eta(w)\) 为 H 的唯一自同构，使得 \(\mathrm{L}(\eta(w)) = \mathrm{L}(w)\)。则映射 \(\eta\) 是从群 Aut G 到群 \(\mathrm{Aut}(\mathrm{H}, \mathrm{N})\) 的同构。
\end{enumerate}

要证明 (i)，根据第 1 号引理 3，只需验证映射 \((u,g)\mapsto\theta(u)g\) 在 \((\mathrm{Id}_{\mathrm{L(H)}},e)\) 的一个邻域内是解析的。存在 L(H) 中 0 的一个开邻域 B，使得 \(\psi = \exp_{\mathbf{H}}|\mathbf{B}\) 是从 B 到 H 中 \(e\) 的一个开邻域的解析同构。存在 Aut L(H) 中 \(\mathrm{Id}_{\mathrm{L(H)}}\) 的一个开邻域 U 和 L(H) 中 0 的一个开邻域 B'，使得 \(\mathrm{U(B')} \subset \mathrm{B}\)。于是映射 \((u, g) \mapsto \theta(u)g\) 从 \(\mathrm{U} \times \psi(\mathrm{B'})\) 到 H 是由以下映射复合而成：

\[
\begin{array}{l l l l} \text {the mapping} (u, g) \mapsto \langle u, \psi^ {- 1} (g) \rangle & \text {of} & \mathrm{U} \times \psi (\mathrm{B} ^ {\prime}) & \text {into} \quad \mathrm{U} \times \mathrm{B} ^ {\prime}; \\ \text {the mapping} (u, x) \mapsto u (x) & \text {of} & \mathrm{U} \times \mathrm{B} ^ {\prime} & \text {into} \quad \mathrm{B}; \\ \text {the mapping} y \mapsto \psi (y) & \text {of} & \mathrm{B} & \text {into} \quad \mathrm{G}. \end{array}
\]

所以这个映射是解析的。

令 \(p\) 为从 \(H\) 到齐性空间 \(H/N\) 的典范映射。则 \(\theta^{-1}(\text{Aut}(H, N))\) 是满足以下条件的 \(u \in \text{Aut } L(H)\) 的集合：

\[
p (\theta (u) g) = p (e), \quad p (\theta (u ^ {- 1}) g) = p (e)
\]

对所有 \(g \in \mathbb{N}\) 均成立。由第 8 节第 2 号定理 2 及其推论 2，这就证明了 (ii)。

假设 N 是离散正规子群。令 \(w \in \operatorname{Aut} G\)。于是

\[
\mathrm{L} (p \circ \eta (w)) = \mathrm{L} (\eta (w)) = \mathrm{L} (w) = \mathrm{L} (w \circ p)
\]

所以 \(p \circ \eta(w) = w \circ p\)，从而 \(\eta(w) \in \text{Aut(H, N)}\)。显然，从 Aut G 到 Aut(H, N) 的映射 \(\eta\) 是单同态。这个同态是满射的，因为 \(p: H \to G\) 是浸没。

设 G 是局部紧群，\(\Gamma\) 是 G 的自同构群。回忆一下，\(T_{\beta}\) 上已经定义了一个拓扑 \(\Gamma\)（《一般拓扑学》第十章第 3 节第 5 号）。它是使映射 \(v \mapsto v\) 和 \(v \mapsto v^{-1}\) 连续的最粗拓扑，其中这些映射从 \(\Gamma\) 到 \(\mathcal{C}_{\varepsilon}(G; G)\)。拓扑 \(T_{\beta}\) 与 \(\Gamma\) 上的群结构相容（同上）。对 G 的每个紧子集 L 和 G 中 \(e_{\mathrm{G}}\) 的每个邻域 U，令 N(L, U) 为满足 \(\phi \in \Gamma\)、\(\phi(g) \in gU\)、\(\phi^{-1}(g) \in gU\) 对所有 \(g \in L\) 成立的集合；则 N(L, U) 构成 \(e_{\Gamma}\) 的邻域基。若 G 由一个紧子集 C 生成，则 \(T_{\beta}\) 也是使映射 \(v \mapsto v[C\) 和 \(v \mapsto v^{-1}[C]\) 从 \(\Gamma\) 到 \(\mathcal{C}_{u}(C; G)\) 连续的最粗拓扑（因为 G 的每个紧子集都包含于 \((\mathrm{C} \cup \mathrm{C}^{-1})^{\mathrm{n}}\)（当 n 足够大时））。若 K 是局部紧的，V 是 K 上的有限维向量空间，则拓扑 \(T_{\beta}\) 正是通常拓扑。

定理 1. 设 G 是有限维李群，\(G_{0}\) 是其单位元分支。假设 G 由 \(G_{0}\) 和有限个元素生成。

\begin{enumerate}
\def\labelenumi{(\roman{enumi})}
\tightlist
\item
  Aut G 上存在唯一的解析流形结构，满足以下条件：
\end{enumerate}

(AUT) 对每个解析流形 M 及每个从 M 到 Aut G 的映射 f，f 是解析的，当且仅当从 M × G 到 G 的映射 \((m, g) \mapsto f(m)g\) 是解析的。

以下均假设 Aut G 具有这一结构。

\begin{enumerate}
\def\labelenumi{(\roman{enumi})}
\setcounter{enumi}{1}
\item
  Aut G 是有限维李群。
\item
  从 Aut G 到 Aut L(G) 的态射 \(\phi: u \mapsto L(u)\) 是解析的。
\item
  若 G 是连通的，则 \(\phi\) 是从李群 Aut G 到 Aut L(G) 的一个李子群的同构；若 G 是单连通的，这个李子群等于 Aut L(G)。
\item
  令 \(\mathfrak{a}\) 为 \(\mathbf{G}\) 的无穷小自同构的集合。则 \(\mathfrak{a}\) 是一个向量场李代数，且与映射 \((u, g) \mapsto u(g)\) 相应的无穷小运算律（该映射从 (Aut G) \(\times\) G 到 G）是从 L(Aut G) 到 \(\mathfrak{a}\) 的同构。
\item
  李群 Aut G 的拓扑就是拓扑 \(\mathcal{T}_{\beta}\)。
\end{enumerate}

\begin{enumerate}
\def\labelenumi{(\alph{enumi})}
\item
  (i) 中所考虑的解析结构的唯一性是显然的。
\item
  假设 G 是连通的。令 H 为 G 的万有覆叠空间，p 为从 H 到 G 的典范态射，且 N = Ker p。~我们引入引理 4 中的记号 \(\theta\)、\(\eta\) 和 \(\operatorname{Aut}(\mathrm{H},\mathrm{N})\)。借助 \(\theta\)，将 Aut L(G) 的李群结构传递给 Aut H。于是 Aut H 成为有限维李群，\(\operatorname{Aut}(\mathrm{H},\mathrm{N})\) 成为 Aut H 的李子群（引理 4 (ii)）。将 \(\operatorname{Aut}(\mathrm{H},\mathrm{N})\) 的李群结构传递给 Aut G，所用映射为 \(\eta^{-1}\)。于是 Aut G 成为有限维李群。定理的性质 (ii)、(iii) 和 (iv) 得到满足，并且映射 \((u,g)\mapsto u(g)\) 从 \((\operatorname{Aut}\mathbf{G})\times\mathbf{G}\) 到 G 是解析的（引理 4 (i)）。设 M 是解析流形，f 是从 M 到 Aut G 的映射，\(\phi\) 是从 M × G 到 G 的映射 \((m,g)\mapsto f(m)g\)。显然，若 f 是解析的，则 \(\phi\) 是解析的。假设 \(\phi\) 是解析的。于是 T \(\phi\) : TM × TG → TG 是解析的；它在 M × L(G) 上的限制，即映射 \((m,x)\mapsto\operatorname{L}(f(m))x\)，因而是解析的；由于 L(G) 是有限维的，可得映射 \(m\mapsto\operatorname{L}(f(m))\) 是解析的，从而 f 是解析的。因此 (i) 成立。
\end{enumerate}

给 \(\mathrm{L}(\mathrm{G})\) 赋予一个范数。对所有 \(\lambda > 0\)，令 \(\mathrm{B}_{\lambda}\) 为以 0 为中心、半径为 \(\lambda\) 的 \(\mathrm{L}(\mathrm{G})\) 中的开球。选择 \(\lambda > 0\) 充分小，使得 \(\psi = \exp_{\mathrm{G}}|\mathrm{B}_{\lambda}\) 是从解析流形 \(\mathrm{B}_{\lambda}\) 到 \(\psi(\mathrm{B}_{\lambda})\) 的 \(G\) 的开子流形的同构。令 \(\Phi\) 是 Aut \(G\) 上的滤子。对于 \(\Phi\)，它收敛到 \(\mathrm{Id}_{\mathrm{G}}\)，且是在 Aut \(G\) 中的收敛，充要条件是 \(\mathrm{L}(\Phi)\) 收敛到 \(\mathrm{Id}_{\mathrm{L}(\mathrm{G})}\)，且是在 Aut \(L(\mathrm{G})\) 中的收敛，从而 \(L(\Phi)|B_{\lambda/2}\) 和 \(L(\Phi)^{-1}\mathrm{B}_{\lambda/2}\) 一致收敛到 \(\mathrm{Id}_{\mathrm{B}_{\lambda/2}}\)。这一条件蕴含 \(\Phi|\psi(\mathrm{B}_{\lambda/2})\) 和 \(\Phi^{-1}|\psi(\mathrm{B}_{\lambda/2})\) 一致收敛到 \(\mathrm{Id}_{\psi(\mathrm{B}_{\lambda/2})}\)。反之，假设 \(\Phi|\psi(\mathrm{B}_{\lambda/2})\) 一致收敛到 \(\mathrm{Id}_{\psi(\mathrm{B}_{\lambda/2})}\)。存在 \(M \in \Phi\)，使得若 \(u \in M\)，则 \(u(\psi(\mathrm{B}_{\lambda/2})) \subset \psi(\mathrm{B}_{2\lambda/3})\)；于是 \(L(u)(\mathrm{B}_{\lambda/2})\) 是 \(L(\mathrm{G})\) 的连通子集，其在 \(\exp_{\mathrm{G}}\) 下的像包含于 \(\psi(\mathrm{B}_{2\lambda/3})\) 中，因此 \(L(u)(\mathrm{B}_{\lambda/2})\) 不与 \(B_{\lambda} - B_{2\lambda/3}\) 相交，从而 \(L(u)(\mathrm{B}_{\lambda/2}) \subset B_{\lambda}\)；于是 \(\Phi|\psi(\mathrm{B}_{\lambda/2})\) 一致收敛到 \(\mathrm{Id}_{\psi(\mathrm{B}_{\lambda/2})}\) 的假设蕴含 \(L(\Phi)|B_{\lambda/2}\) 一致收敛到 \(\mathrm{Id}_{\mathrm{B}_{\lambda/2}}\)。于是得到：

\((\Phi\) 在 Aut G 中收敛到 \(\mathrm{Id}_{\mathrm{G}}\)) \(\Leftrightarrow\)（\(\Phi\) 在 Aut G 中收敛到 \(\mathrm{Id}_{\mathrm{G}}\)，在 \(\mathcal{T}_{\beta}\) 下）。

这就证明了 (vi)。

令 D 为 \(\mathrm{Aut(G)}\) 在 G 上的左作用律所对应的无穷小运算律。由第 1 号命题 1 和 2，\(\mathrm{D(L(AutG))} = a\)。因此 \(a\) 是一个向量场李代数，且 D 是从 \(\mathrm{L(AutG)}\) 到 \(a\) 的态射。设 \(x_{1}\) 和 \(x_{2}\) 是 \(\mathrm{L(AutG)}\) 中满足 \(\mathrm{D}(x_1) = \mathrm{D}(x_2)\) 的元素。则 K 在 G 上的两个运算律 \((\lambda, g) \mapsto (\exp \lambda x_1)g\) 和 \((\lambda, g) \mapsto (\exp \lambda x_2)g\) 具有相同的相应无穷小运算律；因此，当 \(|\lambda|\) 充分小时，\(\exp \lambda x_1\) 和 \(\exp \lambda x_2\) 在 \(e\) 的一个邻域内重合（第 4 节第 7 号定理 6），从而 \(\exp \lambda x_1 = \exp \lambda x_2\)。可得 \(x_1 = x_2\)，所以 D 是从 \(\mathrm{L(AutG)}\) 到 \(a\) 的同构。

因此，定理在 G 连通时已经完全证明。

\begin{enumerate}
\def\labelenumi{(\alph{enumi})}
\setcounter{enumi}{2}
\tightlist
\item
  转到一般情形。根据假设，G 由 \(G_{0}\) 和有限个元素 \(x_{1}, x_{2}, \ldots, x_{n}\) 生成。每个 \(u \in AutG\) 都保持 \(G_{0}\) 稳定。令 \(Aut_{1}G\) 为这样的 \(u \in AutG\) 的集合：它们在取商后诱导出 \(G/G_{0}\) 的恒等自同构。这是 Aut G 的正规子群。由证明的 (b) 部分，Aut \(AutG_{0}\) 具有典范的李群结构，并且映射 \((g_{1}, g_{2}, \ldots, g_{n}, u) \mapsto (ug_{1}, ug_{2}, \ldots, ug_{n})\) 从 \(G_{0}^{n} \times AutG_{0}\) 到 \(G_{0}^{n}\) 是解析的。令 P 为 \(AutG_{0}\) 与 \(G_{0}^{n}\) 相应的半直积；它是有限维李群（第 \(\S\) 1 节第 4 号命题 7）。
\end{enumerate}

若 \(w\in \mathrm{Aut}_1\mathrm{G}\)，记

\[
\begin{array}{r l} w _ {0} & = w | \mathrm{G} _ {0} \in \mathrm{Aut}   \mathrm{G} _ {0} \\ w _ {i} & = x _ {i} ^ {- 1} w (x _ {i}) \in \mathrm{G} _ {0} \quad (1 \leqslant i \leqslant n) \\ \zeta (w) & = ((w _ {1}, \ldots , w _ {n}), w _ {0}) \in \mathrm{P}. \end{array}
\]

对所有 \(w, w'\) 属于 \(\mathrm{Aut}_1\mathrm{G}\)，

\[
\begin{array}{r l} \zeta (w) \zeta (w ^ {\prime}) & = \left((w _ {1}, \ldots , w _ {n}), w _ {0}\right) \left((w _ {1} ^ {\prime}, \ldots , w _ {n} ^ {\prime}), w _ {0} ^ {\prime}\right) \\ & = \left((w _ {1} w _ {0} (w _ {1} ^ {\prime}), \ldots , w _ {n} w _ {0} (w _ {n} ^ {\prime})), w _ {0} w _ {0} ^ {\prime}\right) \\ & = \left((x _ {1} ^ {- 1} w (x _ {1}) w (x _ {1} ^ {- 1} w ^ {\prime} (x _ {1})), \ldots , x _ {n} ^ {- 1} w (x _ {n}) w (x _ {n} ^ {- 1} w ^ {\prime} (x _ {n}))), w _ {0} w _ {0} ^ {\prime}\right) \\ & = \left(((w w ^ {\prime}) _ {1}, \ldots , (w w ^ {\prime}) _ {n}), (w w ^ {\prime}) _ {0}\right) \\ & = \zeta (w w ^ {\prime}) \end{array}
\]

从而 \(\zeta\) 是从 \(\mathrm{Aut}_1\mathrm{G}\) 到 P 的同态。这个同态显然是单射。

我们证明 \(\zeta(\mathrm{Aut}_1\mathrm{G})\) 在 P 中是闭的。令 \(\Phi\) 是 \(\mathrm{Aut}_1\mathrm{G}\) 上的滤子，使得 \(\zeta(\Phi)\) 收敛到 P 中的一点 \(((w_1,\dots,w_n),w_0)\)。于是 \(\Phi\) 逐点收敛到从 G 到 G 的映射 \(v\)。显然，\(v\) 是群 G 的端同态。此外，\(v\) 保持模 \(\mathbf{G}_0\) 的每个陪集稳定，且 \(v |\mathrm{G}_0 = w_0\)。因此 \(v\in \mathrm{Aut}_1\mathrm{G}\)。由于 \(\zeta(v) = ((w_1,\dots,w_n),w_0)\)，我们证明了 \(\zeta(\mathrm{Aut}_1\mathrm{G})\) 在 P 中是闭的。

\begin{enumerate}
\def\labelenumi{(\alph{enumi})}
\setcounter{enumi}{3}
\tightlist
\item
  在证明的 (d) 部分，我们假设 \(\mathbf{K} = \mathbf{R}\)。由第 8 节第 2 号定理 2，\(\zeta(\mathrm{Aut}_1\mathrm{G})\) 是 P 的李子群。将 \(\zeta(\mathrm{Aut}_1\mathrm{G})\) 上的实李群结构传递给 \(\mathrm{Aut}_1\mathrm{G}\)，所用映射为 \(\zeta^{-1}\)。于是 \(\mathrm{Aut}_1\mathrm{G}\) 成为有限维李群。
\end{enumerate}

设 M 是解析流形，f 是从 M 到 \(Aut_{1}G\) 的映射，\(\phi\) 是映射 \((m,g)\mapsto f(m)g\) 从 \(M\times G\) 到 G。有以下等价关系：

\(f\) 是解析的 \(\Leftrightarrow\) 映射 \(m\mapsto (f(m))_i\)，其中 \(0\leqslant i\leqslant n\)，是解析的 \(\Leftrightarrow \left\{ \begin{array}{ll}\text{映射 } m\mapsto f(m)x_i\text{ 从 M 到 G，对 } 1\leqslant i\leqslant n,\text{ 是解析的} \\ \text{且} \\ \text{映射 } (m,g)\mapsto f(m)g\text{ 从 M} \times \mathrm{G}_0\text{ 到 G 是解析的} \end{array} \right.\) \(\Leftrightarrow \phi\) 是解析的。

对于 \(w \in \mathrm{Aut}_1\mathrm{G}\)，有 \(\mathrm{L}(w) = \mathrm{L}(w_0)\)，因此态射 \(w \mapsto \mathrm{L}(w)\) 从 \(\mathrm{Aut}_1\mathrm{G}\) 到 \(\mathrm{Aut}\, \mathrm{L}(\mathrm{G})\) 是解析的。如同 (b) 中一样，\(\mathrm{Aut}_1\mathrm{G}\) 在 \(\mathrm{G}\) 上的作用律所对应的无穷小运算律，是从 \(\mathrm{L}(\mathrm{Aut}_1\mathrm{G})\) 到 \(\mathfrak{a}\) 的同构。

令 C 是 \(G_{0}\) 的一个紧子集，生成 \(G_{0}\)。对于滤子 \(\Phi\)，\(Id_{G}\) 是它在 \(Aut_{1}G\) 上收敛到的极限的充要条件是，\(\Phi|(C \cup \{x_{1}\} \cup \cdots \cup \{x_{n}\})\) 和 \(\Phi^{-1}|(C \cup \{x_{1}\} \cup \cdots \cup \{x_{n}\})\) 一致收敛到

\[
\operatorname{Id} _ {\mathbb {G}} \left| \left(\mathrm{C} \cup \{x _ {1} \} \cup \dots \cup \{x _ {n} \}\right). \right.
\]

因此 \(\mathrm{Aut}_1\mathrm{G}\) 的拓扑就是拓扑 \(\mathcal{T}_{\beta}\)。

显然，\(\mathrm{Aut}_1\mathrm{G}\) 在带有拓扑 \(\mathcal{T}_{\beta}\) 的 Aut G 中是开的。Aut G 上存在与该拓扑相容的李群结构，并在 \(\mathrm{Aut}_1\mathrm{G}\) 上诱导出上述构造的结构（第 8 节第 1 号定理 1 的推论 2）。李群 Aut G 具有定理所述性质这一事实，由 \(\mathrm{Aut}_1\mathrm{G}\) 的相应性质推出。

\begin{enumerate}
\def\labelenumi{(\alph{enumi})}
\setcounter{enumi}{4}
\tightlist
\item
  在证明的 (e) 部分，我们假设 \(\mathbf{K} = \mathbf{C}\)。由 (c) 以及第 8 节第 2 号定理 2，\(\mathrm{Aut}_1\mathrm{G}\) 上存在实李群结构，使得 \(\zeta\) 是从 \(\mathrm{Aut}_1\mathrm{G}\) 到 \(\mathbb{P}\) 的一个实李子群的同构。
\end{enumerate}

\((w, g) \mapsto wg\) 在 \((\mathrm{Aut}_1\mathrm{G}) \times \mathrm{G}\) 上、即在 \(\mathbf{G}\) 上的运算律是实解析的。令 \(\mathbf{D}\) 为相应的无穷小运算律。由第 1 号命题 1 和 2，\(\mathrm{D}(\mathrm{L}(\mathrm{Aut}_1\mathrm{G})) = \mathfrak{a}\)。

对所有 \(\alpha \in \mathfrak{a}\)，令 \(\alpha_0\) 表示 \(\alpha\) 在 \(\mathrm{G}_0\) 上的限制；它是 \(\mathrm{G}_0\) 的无穷小自同构，由于证明的 (b) 部分，我们将它与 \(\mathrm{L}(\mathrm{Aut}\mathrm{G}_0)\) 中的一个元素识别。对 \(1 \leqslant i \leqslant n\)，记

\[
\alpha_ {i} = x _ {i} ^ {- 1} \alpha (x _ {i}) \in \mathrm{L} (\mathrm{G}) = \mathrm{L} (\mathrm{G} _ {0}).
\]

最后，记 \(f(\alpha) = ((\alpha_1, \ldots, \alpha_n), \alpha_0) \in \mathrm{L}(\mathrm{P})\)。则 \(f\) 是一个 \(\mathbf{C}\)-线性映射，作用于 \(\mathfrak{a}\) 并取值于 \(\mathrm{L}(\mathrm{P})\)。

另一方面，显然 \(\mathrm{L}(\zeta) = f\circ \mathrm{D}\)。因此 \(\mathrm{L}(\zeta)(\mathrm{L}(\mathrm{Aut}_1\mathrm{G})) = f(a)\) 是 \(\mathrm{L}(P)\) 的复向量子空间。由第 4 节第 2 号命题 2，\(\zeta (\mathrm{Aut}_1\mathrm{G})\) 是 \(P\) 的复李子群，我们可以完全按照 (d) 中的方式进行：将 \(\zeta (\mathrm{Aut}_1\mathrm{G})\) 上的复李群结构传递给 \(\mathrm{Aut}_1\mathrm{G}\)，所用映射为 \(\zeta^{-1}\)，并如 (d) 中那样看出，\(\mathrm{Aut}_1\mathrm{G}\) 具有与定理的性质 (i)、(ii)、(iii)、(v) 和 (vi) 相应的性质。

显然 \(\mathrm{Aut}_1\mathrm{G}\) 在 Aut G 中带有拓扑 \(\mathcal{T}_{\beta}\) 时是开的。令 \(w\in \mathrm{Aut}G\)。令 \(\sigma\) 为自同构 \(v\mapsto wvw^{-1}\)，作用在 \(\mathrm{Aut}_1\mathrm{G}\) 上。它是实解析的，\(\mathrm{L}(\sigma)\) 是一个 \(\mathbf{R}\)-自同构 \(\mathrm{L}(\mathrm{Aut}_1\mathrm{G})\)，并且

\[
\mathrm{D} \circ \mathrm{L} (\mathrm{Aut} _ {1} \mathrm{G}) \circ \mathrm{D} ^ {- 1}
\]

是 a 的一个 R-自同构。这个自同构也是由 w 通过结构传递得到的 a 的自同构；由于 w 是 K-解析的，我们看到 \(L(\sigma)\) 是 K-线性的。因此 \(\sigma\) 是 K-解析的（第 \(\S\) 3 节第 8 号命题 32）。由第 \(\S\) 1 节第 9 号命题 18，Aut G 上存在唯一的李 K-群结构，使得 \(Aut_1G\) 是 Aut G 的开李子群。这一结构具有定理所述性质这一事实，由 \(Aut_1G\) 的相应性质推出。

推论 1. 设 G 是有限维实李群，\(G_{0}\) 是其单位元分支。假设 G 由 \(G_{0}\) 和有限个元素生成。则 Aut G 具有拓扑 \(T_{\beta}\)，并且是有限维实李群。

推论 2. 设 G 是半单连通实或复李群。群 Int G 是 Aut G 的单位元分支。

映射 \(u \mapsto \mathrm{L}(u)\) 是从 Aut G 到 Aut L(G) 的一个李子群的同构（定理 1）。Int G 在此同构下的像是 Ad G。但 Ad G 是 Aut L(G) 的单位元分支（第 9 节第 8 号命题 30(ii)）。

\subsection{李群的自同构群（超度量情形）}

定理 2. 当 K 是超度量且局部紧的、G 是紧李群时，定理 1 的断言 (i)、(ii)、(iii)、(v) 和 (vi) 成立。

\begin{enumerate}
\def\labelenumi{(\alph{enumi})}
\item
  (i) 中所考虑的解析结构的唯一性是显然的。
\item
  假设 G 是由可赋范李代数 L 定义的李群。则 G 是 L 中的一个开闭球。令 \(w \in \operatorname{Aut} G\)。于是 \(\operatorname{L}(w)\) 在 0 的一个邻域内与 w 重合。令 \(x \in G\)。令 p 为剩余域的特征。于是 \(p^{n}x\) 当 n 趋于 \(+\infty\) 时趋于 0。因此存在 n，使得 \(w(p^{n}x) = \operatorname{L}(w)(p^{n}x)\)。所以
\end{enumerate}

\[
\begin{array}{r l} p ^ {n} w (x) & = w (x) ^ {p ^ {n}} = w (x ^ {p ^ {n}}) = w (p ^ {n} x) \\ & = \mathrm{L} (w) (p ^ {n} x) = p ^ {n} \mathrm{L} (w) (x) \end{array}
\]

从而 \(w(x) = \mathrm{L}(w)(x)\)。因此，\(w = \mathrm{L}(w)|\mathrm{G}\)。

令 \(\Gamma\) 为这样的 \(\gamma \in \operatorname{Aut} L(G)\) 的集合，满足 \(\gamma(G) = G\)。由于 \(G\) 在 \(L(G)\) 中既开又紧，\(\Gamma\) 是 \(\operatorname{Aut} L(G)\) 的开子群。由上面所述，Aut \(G\) 可与 \(\Gamma\) 识别，从而 Aut \(G\) 上存在李群结构，在此结构下定理 1 的性质 (i)、(ii)、(iii) 和 (vi) 显然成立。性质 (v) 由第 1 号命题 1 和 3 得出。

\begin{enumerate}
\def\labelenumi{(\alph{enumi})}
\setcounter{enumi}{2}
\tightlist
\item
  转到一般情形。由第 7 节第 1 号命题 1，存在一个开紧子群 \(G_0\)，它是 \(G\) 的子群，属于 (b) 中所考虑的类型。于是 \(G\) 由 \(G_0\) 和有限个元素 \(x_1, x_2, \ldots, x_n\) 生成。令 \(\text{Aut}_1G\) 为这样的 \(u \in \text{Aut}G\) 的集合：满足 \(u(G_0) = G_0\) 且 \(u(x_iG_0) = x_iG_0\)（\(1 \leqslant i \leqslant n\)）。按照定理 1 的证明 (c) 部分，定义 \(P\)，即 \(\text{Aut}G_0\) 与 \(G_0^n\) 的半直积，以及单同态 \(\zeta\)，它从 \(\text{Aut}_1G\) 到 \(P\)，其像在 \(P\) 中是闭的。
\end{enumerate}

(d)、(e)：论证完全同定理 1 证明的 (d)、(e) 部分，只需将 \(\mathbf{R}\) 换成 \(\mathbf{Q}_p\)，并用命题 3 代替命题 2。

注。若 \(\mathbf{K} = \mathbf{Q}_p\) 且李群 \(\mathbf{G}\) 由一个紧子集生成（参见~练习 2），则定理 1 的断言 (i)、(ii)、(iii) 和 (vi) 仍然成立，但 (v) 不成立（练习 3）。

\appendixsection{线性表示上的运算}


设 \(G\) 是一个群，\(k\) 是交换域，\(E_1, E_2, \ldots, E_n\) 是 \(k\) 上的向量空间，\(\pi_i\) 是 \(G\) 在 \(E_i\) 上的线性表示（\(1 \leqslant i \leqslant n\)）。映射

\[
g \mapsto \pi_ {1} (g) \otimes \dots \otimes \pi_ {n} (g)
\]

是从 G 到向量空间 \(E_{1} \otimes \cdots \otimes E_{n}\) 的线性映射，称为 \(\pi_{1}, \ldots, \pi_{n}\) 的张量积，记作 \(\pi_{1} \otimes \cdots \otimes \pi_{n}\)。

设 E 是 k 上的向量空间，\(\pi\) 是 G 在 E 上的线性表示。对所有 \(g \in G\)，令 \(\tau(g)\)（分别地，\(\sigma(g)\)、\(\varepsilon(g)\)）为代数 \(\mathsf{T}(\mathsf{E})\)（分别地，\(\mathsf{S}(\mathsf{E})\)、\(\bigwedge(\mathsf{E})\)）中延拓 \(\pi(g)\) 的唯一自同构（《代数》第三章第 5 节第 2 号、第 6 节第 2 号和第 7 节第 2 号）。于是 \(\tau\)（分别地，\(\sigma\)、\(\varepsilon\)）是 G 在 T(E)（分别地，\(\mathsf{S}(\mathsf{E})\)、\(\bigwedge(\mathsf{E})\)）上的线性表示，记作 \(\mathsf{T}(\pi)\)（分别地，\(\mathsf{S}(\pi)\)、\(\bigwedge(\pi)\)）。由 T(π)（分别地，\(\mathsf{S}(\pi)\)、\(\bigwedge(\pi)\)）定义于 \(\mathsf{T}^{\mathsf{n}}(\mathsf{E})\)（分别地，\(\mathsf{S}^{\mathsf{n}}(\mathsf{E})\)、\(\bigwedge^{\mathsf{n}}(\mathsf{E})\)）的子表示，称为 \(\pi\) 的 n 次张量（分别地，对称、外）幂，记作 \(\mathsf{T}^{\mathsf{n}}(\pi)\)（分别地，\(\mathsf{S}^{\mathsf{n}}(\pi)\)、\(\bigwedge^{\mathsf{n}}(\pi)\)）。于是

\[
T ^ {n} (\pi) = \pi \otimes \pi \otimes \dots \otimes \pi
\]

（n 个因子）。表示 \(\mathsf{S}(\pi), \bigwedge (\pi)\) 是 \(\mathsf{T}(\pi)\) 的商表示，因此 \(\mathsf{S}^n (\pi), \bigwedge^n (\pi)\) 是 \(\mathsf{T}^n (\pi)\) 的商表示。

设 g 是 k 上的李代数。有限个 g 的表示的张量积已在第一章第 3 节第 2 号定义；记作 \(\pi_{1}\otimes\cdots\otimes\pi_{n}\)。设 E 是 k 上的向量空间，\(\pi\) 是 g 在 E 上的表示。对所有 \(x\in g\)，令 \(\tau'(x)\)（分别地，\(\sigma'(x)\)、\(\varepsilon'(x)\)）为代数 \(\mathrm{T(E)}\)（分别地，\(\mathrm{S(E)}\)、\(\bigwedge(\mathrm{E})\)）中延拓 \(\pi(x)\) 的唯一导子（《代数》第三章第 10 节第 9 号例 1）。于是由《代数》同上，公式 (35)，\(\tau'\)（分别地，\(\sigma'\)、\(\varepsilon'\)）是 g 在 T(E)（分别地，S(E)、\(\bigwedge(\mathrm{E})\)）上的线性表示，

记作 \(\mathsf{T}(\pi)\)（分别地，\(\mathsf{S}(\pi),\bigwedge (\pi)\)）。由 \(\mathsf{T}(\pi)\)（分别地，\(\mathsf{S}(\pi),\bigwedge (\pi)\)）定义的、由 \(\mathsf{T}^n (\mathsf{E})\)（分别地，\(\mathsf{S}^n (\mathsf{E}),\bigwedge^n (\mathsf{E}))\)）决定的子表示，记作 \(\mathsf{T}^n (\pi)\)（分别地，\(\mathsf{S}^n (\pi),\bigwedge^n (\pi))\)）。表示 \(\mathsf{T}^n (\pi)\) 是 \(n\) 个与 \(\pi\) 相同的表示的张量积。表示 \(\mathsf{S}(\pi),\bigwedge (\pi)\) 是 \(\mathsf{T}(\pi)\) 的商表示，因此 \(\mathsf{S}^n (\pi),\bigwedge^n (\pi)\) 是 \(\mathsf{T}^n (\pi)\) 的商表示。

\specialchapter{习题}

\exercisesfor{1}

\begin{enumerate}
\def\labelenumi{\arabic{enumi}.}
\tightlist
\item
  设 G 是有限维连通实或复李群。设 H 和 L 是 G 的李子群，满足 HL = G。证明典范映射 \(G \rightarrow G/H\) 和 \(G \rightarrow G/L\) 定义出解析流形的一个同构
\end{enumerate}

\[
\mathrm{G} / (\mathrm{H} \cap \mathrm{L}) \rightarrow (\mathrm{G} / \mathrm{H}) \times (\mathrm{G} / \mathrm{L}).
\]

\begin{enumerate}
\def\labelenumi{\arabic{enumi}.}
\setcounter{enumi}{1}
\item
  设 \(\mathrm{P} = \{z\in \mathbf{C}|\mathcal{I}(z) > 0\}\)。设 \(\mathbf{G}\) 是由双射 \(z\mapsto az + b\) 在 \(\mathrm{P}(a > 0,b\in \mathbf{R})\) 上组成的群。于是 \(\mathbf{G}\) 在 \(\mathrm{P}\) 上单纯传递地作用，这使我们能够将 \(\mathbf{P}\) 的复解析流形结构传递给 \(\mathbf{G}\)。证明所得结构在 \(\mathbf{G}\) 的左平移下不变，但在右平移下不变。
\item
  设 \(R_{d}\) 是带离散流形结构的实数加法群。令 \(R_{d}\) 按运算律 \((x,y)\mapsto x+y\) 作用在解析流形 R 上。则 \(R_{d}\) 在 R 上传递地作用，但 R 不是 \(R_{d}\) 的李齐性空间。
\item
  设 H 是一个紧的实李群，作用于有限维实流形 V；设 V 属于类 \(C^{r}\)，其中 \(r \in N_{R}\)，且 H 在 V 上的作用属于类 \(C^{r}\)。设 \(p \in V\) 在 H 下不变；H 在线性地作用于 V 在 p 处的切空间 T 上。~证明存在 p 的一个 H-稳定开邻域 U 和一个 \(C^{r}\) -态射 \(f: U \to T\)，使得：
\end{enumerate}

\begin{enumerate}
\def\labelenumi{(\alph{enumi})}
\item
  \(f(p) = 0\)，且 \(f\) 在 \(p\) 处的切映射是恒等映射；
\item
  \(f\) 与 \(\mathbf{H}\) 在 \(\mathbf{U}\) 和 \(\mathbf{T}\) 上的作用相交换。
\end{enumerate}

(先选取 \(f_0\)，使其满足 \((a)\)，再按公式定义 \(f\)

\[
f (x) = \int_ {\mathbf {H}} h. f _ {0} (h ^ {- 1} x) d h,
\]

其中 \(dh\) 是 \(H\) 上总质量为 1 的 Haar 测度。)

推出 H-空间 V 和 T 在 p 的某个邻域内局部同构；换言之，存在 V 在 p 处的坐标系，使得 H 关于这些坐标线性作用（``Bochner 定理``）。

\begin{enumerate}
\def\labelenumi{\arabic{enumi}.}
\setcounter{enumi}{4}
\tightlist
\item
  对超度量域 K 上的流形完成练习 4，假设 H 是一个有限群，其阶 n 满足 \(n.1 \neq 0\) 在 K 中不为 0。
\end{enumerate}

¶ 6. 设 G 是一个实李群，适当地作用于有限维 Hausdorff 实解析流形 V。设 \(p \in V\)，令 H 为 \(p\) 的稳定子群，令 Gp 为 \(p\) 的轨道。H 是紧的，且流形 Gp 可与李齐性空间 G/H 识别。

\begin{enumerate}
\def\labelenumi{(\alph{enumi})}
\item
  证明存在一个穿过 p 的 V 的子流形 S，在 H 下稳定，并且 \(\mathrm{T}_{p}(\mathrm{V})\) 是 \(\mathrm{T}_{p}(\mathrm{G}p)\) 与 \(\mathrm{T}_{p}(\mathrm{S})\) 的直和。（选取 \(\mathrm{T}_{p}(\mathrm{G}p)\) 在 \(\mathrm{T}_{p}(\mathrm{V})\) 中的一个在 H 下稳定的补空间，并应用练习 4。）
\item
  假设 S 如上选取。H 按 \(h.(g,s)=(gh,h^{-1}s)\) 适当地且自由地作用于 G × S；令 \(\mathrm{E}=(\mathrm{G}\times\mathrm{S})/\mathrm{H}\) 为相应的商流形（参见~《微分与解析流形》，R，6.5.1）。映射 \((g,s)\mapsto gs\) 在取商后定义一个态射 \(\rho:E\to V\)。证明 \(\rho\) 与 G 在 E 和 V 上的作用相交换，并且存在一个包含 p、在 H 下稳定的 S 的开子流形 \(S'\)，使得 \(\rho\) 将 \(\mathrm{E}'=(\mathrm{G}\times\mathrm{S}')/\mathrm{H}\) 诱导为轨道 Gp 的一个饱和邻域上的流形同构。
\item
  令 \(N_{p} = T_{p}(V)/T_{p}(Gp)\) 为在 p 处横截于 Gp 的横截空间（《微分与解析流形》，R，5.8.8）；H 作用于 \(N_{p}\)。由 (b) 推出，知道 H 及 H 在 \(N_{p}\) 上的线性表示，就能在 p 的轨道的某个邻域内确定 G-空间 V。
\item
  若 \(x \in \mathbf{N}_p\)，令 \(\mathbf{H}_x\) 为 \(x\) 在 \(\mathbf{H}\) 中的稳定子群。证明存在一个饱和邻域 \(V'\)，它是 \(Gp\) 的邻域，具有如下性质：
\end{enumerate}

对所有 \(p' \in V'\)，存在 \(x \in \mathbb{N}_p\)，使得 \(p'\) 在 \(G\) 中的稳定子群与（在 \(G\) 中）\(H_x\) 共轭。

\begin{enumerate}
\def\labelenumi{(\alph{enumi})}
\setcounter{enumi}{4}
\tightlist
\item
  证明 \(H_{x}\) 在 H 中按共轭仅有有限多个。（对 \(N_{p}=n\) 作归纳，利用 H 在一个由 H 稳定的 n-1 维球面上的作用。）†
\end{enumerate}

\begin{enumerate}
\def\labelenumi{\arabic{enumi}.}
\setcounter{enumi}{6}
\tightlist
\item
  设 E 是 R 上的完备可赋范空间，F 是 E 的闭向量子空间，并且不存在从 E 到 \(\mathbf{F} \times (\mathbf{E}/\mathbf{F})\) 的双连续线性双射。令 X 为解析流形 \(R \times E \times F\)。令 G 为 F 的加法群，它通过映射
\end{enumerate}

\[
(g, (\lambda , e, f)) \mapsto (\lambda , e + g, f + \lambda g)
\]

从 \(\mathbf{G} \times \mathbf{X}\) 到 \(\mathbf{X}\) 的作用映射。对所有 \(x \in \mathbf{X}\)，令 \(\mathbf{R}_x\) 为 \(x\) 的 G-轨道，它是 \(\mathbf{X}\) 的准子流形。于是命题 10 的条件 (a) 得到满足，但不存在具有下列性质的有序对 \((\mathbf{Y}, \pi)\)：\(\mathbf{Y}\) 是解析流形，\(\pi\) 是从 \(\mathbf{X}\) 到 \(\mathbf{Y}\) 的态射，且对所有 \(x \in \mathbf{X}\)，

\[
\mathrm{T} _ {x} (\pi): \mathrm{T} _ {x} (\mathrm{X}) \rightarrow \mathrm{T} _ {\pi (x)} (\mathrm{Y})
\]

是满射的，且核为 \(\mathrm{T}_x(\mathbf{R}_x)\)。

（假设 \((\mathrm{Y},\pi)\) 存在。令 \(\mathbf{H}\) 为 \(\mathbf{Y}\) 在 \(\pi (0,0,0)\) 处的切空间。于是 \(\mathbf{H}\) 同构于 \(\mathbf{R}\times \mathbf{F}\times (\mathrm{E} / \mathrm{F})\)。另一方面，\(\mathbf{H}\) 与 \(\mathrm{T}_{\pi (x)}\mathrm{Y}\) 同构，当 \(x\) 充分接近 \((0,0,0)\) 时，因而同构于 \(\mathbf{R}\times \mathbf{E}\)。）

\begin{enumerate}
\def\labelenumi{\arabic{enumi}.}
\setcounter{enumi}{7}
\item
  令 \(\mathbf{T}\) 带有归一化 Haar 测度，并在 \(\mathcal{L}(\mathrm{L}^2 (\mathbf{T}))\) 上赋予范数拓扑。证明 \(\mathbf{T}\) 在 \(\mathrm{L}^2 (\mathbf{T})\) 上的正则表示不连续。（对所有 \(g\in \mathbf{T}\)（其不同于 \(\epsilon\)），构造 \(f\in \mathrm{L}^2 (\mathbf{T})\)，使得 \(\| f\| = 1\)，\(\| \gamma (g)f - f\| = \sqrt{2}\)。）
\item
  令 I 是由 \((x, \sqrt{2} x) \in \mathbf{R}^2\) 组成的集合，其中 \(-\frac{1}{2} < x < \frac{1}{2}\)。令 U 为 I 在实李群 \(\mathbf{H} = \mathbf{T}^2\) 中的典范像。证明 U 是 H 的李子群芽，但由 U 生成的子群 \(\mathrm{H}'\) 在 H 中稠密且不闭。
\item
  令 \(\mathbf{R}_d\) 为带离散流形结构的群 \(\mathbf{R}\)。令 \(\mathrm{H}\) 为实李群 \(\mathbf{R} \times \mathbf{R}_d\)。令 \(\mathrm{U}\) 为 \(\mathrm{H}\) 中形如 \((x, x)\) 或形如 \((x, -x)\) 的元素集合，其中 \(x \in \mathbf{R}\)。于是 \(\mathrm{U}\) 是 \(\mathrm{H}\) 的李子群芽，\(\mathrm{U}\) 是离散的，\(\mathrm{G}\) 是 \(\mathrm{H}\) 的由 \(\mathrm{U}\) 生成的子群，且等于 \(\mathrm{H}\)；命题 22 的推论在 \(\mathrm{G}\) 上定义的李群结构是离散群结构。
\end{enumerate}

\exercisesfor{3}

\begin{enumerate}
\def\labelenumi{\arabic{enumi}.}
\item
  设 \(G\) 是一个李群，\(G_1\) 和 \(G_2\) 是 \(G\) 的李子群。假设 \(G_2\) 在 \(G\) 中具有有限余维，且 \(K\) 的特征为 0。则 \(G_1 \cap G_2\) 是 \(G\) 的李子群，其李代数为 \(L(G_1) \cap L(G_2)\)（仿照命题 29 及其推论 2 论证）。
\item
  假设 \(\mathbf{K}\) 是局部紧的。令 \(\mathbf{G}\) 是维数为 \(n\) 的李群。
\end{enumerate}

\begin{enumerate}
\def\labelenumi{(\alph{enumi})}
\tightlist
\item
  设 \(\phi\) 是 G 的具有有限核的 étale 端同态。令 \(\mu\) 是 G 的左 Haar 测度。则 \(\phi\) 是适当的，且
\end{enumerate}

\[
\phi (\mu) = \alpha \cdot \mu | \phi (G) \quad \text { where } \quad \alpha = \frac {\operatorname{Card} (\operatorname{Ker} \phi)}{\operatorname{mod} \det L (\phi)}.
\]

(使用命题 55。)
